\documentclass[11pt]{article}

% ============================================================================
% Closure axis (Paper B, flagship) — modular statements-of-record root.
%
% main.tex carries the preamble + front matter + the ordered \input of the
% sections/ and appendices/ files (the labeled environments live there; the
% inventory extractor reads those files). Preamble adopted from
% ../paper-template (publication style: lmodern/microtype, titling, fancyhdr
% footer, orcidlink author, caption/booktabs); BSD macros are in
% includes/paper_macros.tex.
%
% Discipline: this is the CONDITIONAL Strong-BSD closure — NOT an
% unconditional proof. The deep arithmetic enters only as imported
% proof-carrying structures and typed recognition-source carriers; never an
% axiom, mathlib-free. Label schema <kind>:closure:<short-name>.
% ============================================================================

\usepackage[T1]{fontenc}
\usepackage[utf8]{inputenc}
\usepackage[english]{babel}
\usepackage[margin=1in]{geometry}
\usepackage{lmodern}
\usepackage{microtype}
\usepackage{mathtools}
\usepackage{amsmath,amssymb,amsthm}
\usepackage{array}
\usepackage{booktabs}
\usepackage{tabularx}
\usepackage{longtable}
\usepackage{float}
\usepackage{graphicx}
\usepackage{tikz}
\usetikzlibrary{arrows.meta,calc,fit,positioning}
\usepackage{pdflscape}
\usepackage{caption}
\usepackage{enumitem}
\usepackage{fancyhdr}
\usepackage[numbers]{natbib}
\usepackage{doi}
\usepackage{xurl}
\usepackage{hyperref}
\usepackage{cleveref}
\usepackage{orcidlink}
\usepackage{titling}

\usepackage{aliascnt}
\newtheorem{theorem}{Theorem}[section]
% Shared-counter environments get alias counters so that cleveref prints the
% environment name (Definition, Obligation, ...) rather than "Theorem".
\newcommand{\sharedtheorem}[2]{%
  \newaliascnt{#1}{theorem}%
  \newtheorem{#1}[#1]{#2}%
  \aliascntresetthe{#1}}
\sharedtheorem{lemma}{Lemma}
\sharedtheorem{proposition}{Proposition}
\sharedtheorem{corollary}{Corollary}
\theoremstyle{definition}
\sharedtheorem{definition}{Definition}
\sharedtheorem{example}{Example}
\sharedtheorem{remark}{Remark}
\sharedtheorem{warning}{Warning}
\sharedtheorem{obligation}{Obligation}
\sharedtheorem{nonclaim}{Nonclaim}

\urlstyle{same}
\captionsetup{font=small, labelfont=bf, labelsep=period}
\captionsetup[table]{skip=6pt, position=top}
\captionsetup[figure]{skip=6pt}
\crefname{theorem}{theorem}{theorems}
\Crefname{theorem}{Theorem}{Theorems}
\crefname{lemma}{lemma}{lemmas}
\Crefname{lemma}{Lemma}{Lemmas}
\crefname{proposition}{proposition}{propositions}
\Crefname{proposition}{Proposition}{Propositions}
\crefname{corollary}{corollary}{corollaries}
\Crefname{corollary}{Corollary}{Corollaries}
\crefname{definition}{definition}{definitions}
\Crefname{definition}{Definition}{Definitions}
\crefname{remark}{remark}{remarks}
\Crefname{remark}{Remark}{Remarks}
\crefname{warning}{warning}{warnings}
\Crefname{warning}{Warning}{Warnings}
\crefname{obligation}{obligation}{obligations}
\Crefname{obligation}{Obligation}{Obligations}
\crefname{nonclaim}{nonclaim}{nonclaims}
\Crefname{nonclaim}{Nonclaim}{Nonclaims}
\crefname{figure}{Figure}{Figures}
\Crefname{figure}{Figure}{Figures}
\crefname{table}{Table}{Tables}
\Crefname{table}{Table}{Tables}
\crefname{section}{Section}{Sections}
\Crefname{section}{Section}{Sections}
\crefname{example}{example}{examples}
\Crefname{example}{Example}{Examples}
\crefname{appendix}{appendix}{appendices}
\Crefname{appendix}{Appendix}{Appendices}
\setlist[enumerate]{leftmargin=*, itemsep=2pt, topsep=4pt}
\setlength{\droptitle}{-3.2em}
\pretitle{\begin{center}\LARGE\bfseries}
\posttitle{\par\end{center}\vspace{0.5em}}
\preauthor{\begin{center}\normalsize}
\postauthor{\par\end{center}\vspace{-0.4em}}
\predate{\begin{center}\small}
\postdate{\par\end{center}\vspace{0.6em}}
\clubpenalty=10000
\widowpenalty=10000
\displaywidowpenalty=10000
\interfootnotelinepenalty=10000
\allowdisplaybreaks[2]
\fancypagestyle{firstpage}{\fancyhf{}
  \fancyfoot[C]{\scriptsize Automorph Inc., Wilmington, DE, USA \quad \texttt{ioannis@automorph.io} \quad \textcopyright\ 2026 \quad CC-BY 4.0 \quad \doi{10.5281/zenodo.23130788}}
  \renewcommand{\headrulewidth}{0pt}
  \renewcommand{\footrulewidth}{0pt}
}

% BSD macros + the math-heavy-skeleton layout settings (\emergencystretch, \sloppy).
% Paper macros for the closure axis (Phase 1). \input by main.tex.
% Reserved-symbol policy binds Lean naming (see paper/notation_and_terminology.md).
% This file is the canonical macro source; paper/notes/macro-audit.md mirrors it.

% --- Shared / framework macros (identical block in both papers) ---
\newcommand{\Q}{\mathbb{Q}}
\newcommand{\R}{\mathbb{R}}
\newcommand{\Z}{\mathbb{Z}}
\newcommand{\Qp}{\mathbb{Q}_p}
\newcommand{\Zp}{\mathbb{Z}_p}
\newcommand{\Sha}{\operatorname{Sha}}
\newcommand{\Reg}{\operatorname{Reg}}
\newcommand{\Tam}{\operatorname{Tam}}
\newcommand{\tors}{\operatorname{tors}}
\newcommand{\tr}{\operatorname{tr}}
\newcommand{\Pf}{\operatorname{Pf}}
\newcommand{\Fitt}{\operatorname{Fitt}}
\newcommand{\cores}{\operatorname{cores}}
\newcommand{\rank}{\operatorname{rank}}
\newcommand{\Hf}{H^1_f}
\newcommand{\kapr}{\kappa_r}

% --- Closure-specific macros ---
\newcommand{\Sel}{\operatorname{Sel}}
\newcommand{\SelBSD}{\mathrm{Sel}^{!}_{\mathrm{BSD}}}   % the trace shell
\newcommand{\rhoE}{\rho_E}                              % shell residual
\newcommand{\PiBSD}{\Pi_{\mathrm{BSD}}}                 % recognition predicate
\newcommand{\GammaPD}{\Gamma_{\mathrm{BSD}}^{\mathrm{padic\text{-}descent}}}
\newcommand{\GammaSP}{\Gamma_{\mathrm{BSD}}^{\mathrm{Sha\text{-}persistence}}}
\newcommand{\GammaGZ}{\Gamma_{\mathrm{BSD}}^{\mathrm{higher\text{-}GZ\text{-}fixity}}}
\newcommand{\chiCTp}{\chi_{CT,p}}                       % Cassels--Tate comparison
\newcommand{\AofE}{\mathcal{A}_E}                        % Selmer-complex pairing substrate (calligraphic, distinct from apparatus A_E)
\newcommand{\etap}{\eta_p}                              % OC eta
\newcommand{\RefStable}{\mathbf{RefStable\,AOR}}        % AOR membership

% Math-heavy SoR skeleton: allow interword stretch to avoid overfull boxes
\setlength{\emergencystretch}{3.5em}
\sloppy


\title{A Conditional Closure of the Strong Birch--Swinnerton-Dyer Conjecture}
\author{Ioannis Tsiokos\,\orcidlink{0009-0009-7659-5964}}
\date{v1: 31 August 2026\quad\textperiodcentered\quad v2: 4 October 2026}
\hypersetup{
  pdftitle={A Conditional Closure of the Strong Birch--Swinnerton-Dyer Conjecture},
  pdfauthor={Ioannis Tsiokos},
  hidelinks
}

\begin{document}
\maketitle
\thispagestyle{firstpage}

\begin{abstract}
We formulate the Strong Birch--Swinnerton-Dyer formula for an elliptic
curve $E/\Q$ as a fixed-point statement and prove a conditional theorem
in that formulation.  Let $L^*$ be the leading Taylor coefficient of
$L(E,s)$ at $s=1$ and $R_E$ the arithmetic side of the formula.  We show
that the Strong BSD identity holds exactly when the pair $(L^*,R_E)$ is
fixed by the averaging closure $c(a,b)=((a+b)/2,(a+b)/2)$, and we
determine what a closure on a richer layer of arithmetic data would need
in order to return this statement: a readout intertwining it with $c$,
and fixedness of the original state.  The main theorem derives the
identity from a recognition hypothesis, which in rank at most one
contains the identity and in rank at least two is a leading-term
identity with the normalization
$\kappa_r=\#\Sha/\#E(\Q)_{\mathrm{tors}}^2$.  This hypothesis has the
strength of the conclusion, so the theorem is a conditional translation
of Strong BSD into the closure framework, not an independent derivation.
The other inputs, local descent and signed hypotheses and comparison
statements from the literature, are carried with their precise scope
but are not used for the scalar identity.  We show by examples what these
inputs contain beyond coarser data: for $121a1$ and $121b1$ the local
congruences cannot see the Tamagawa numbers $1$ and $2$, and for $1913b1$
at $p=3$ two generators of the same Selmer-complex determinant ideal
differ in whether they have a square root.  We also show that sensitivity
of the BSD product to its factors does not make any hypothesis necessary,
and we record a proposed overconvergent $\eta$-formula only as a
conditional certificate, explaining why the published results it was
modelled on do not apply as stated.  The paper does not prove BSD.  The
scalar algebra and the closure statements are proved in Lean~4; the
remaining statements are labelled by the exact extent of their
formalization.
\end{abstract}

\medskip
\noindent\textbf{Keywords.} Birch--Swinnerton-Dyer conjecture; Strong BSD;
closure operator; conditional theorem; Cassels--Tate pairing; Selmer
complex; Iwasawa main conjecture; formal verification (Lean); Six Birds
Theory; emergence calculus.

\medskip
\noindent
\textbf{Conventions.} $E/\Q$ is an elliptic curve, $M_E=h^1(E)(1)$,
$r=\operatorname{ord}_{s=1}L(E,s)$ is the analytic rank, and curve labels
are Cremona labels.  The normalization of $\kapr$ is taken from
\citet{TsiokosBSDApparatus2026}.

% ===========================================================================
\section{Introduction}\label{sec:introduction}
% statements-of-record rows resolved here: none

What would a conditional closure of Strong Birch--Swinnerton-Dyer~\citep{BirchSwinnertonDyer1965,Tate1966bsd} look
like, and what exactly would it rest on?  The classical landscape already
contains powerful conditional and partial routes: Gross--Zagier and
Kolyvagin~\citep{GrossZagier1986,Kolyvagin1989} control the rank-one Heegner setting; Skinner--Urban and
Iwasawa-main-conjecture methods~\citep{SkinnerUrban2014} control large ordinary regions;
Cassels--Tate and Selmer-complex formalisms~\citep{Flach1990,Nekovar2006} control part of the
quadratic and determinant bookkeeping.  These results do not, by
themselves, give a single uniform proof of the strong formula in every
rank and at every bad prime.  The question addressed here is narrower
and more structural: if the remaining arithmetic inputs are isolated
with precise names, can one assemble them into a single audited closure
whose scalar readout is exactly the Strong-BSD identity?

The answer is deliberately conditional.  This is \emph{not} an
unconditional classical proof of BSD, scalar or strong.  The result is
theorem-grade within the Six Birds closure discipline, conditional on
the standing shell-closure record, the named recognition sources, and the
headline imported theorem stacks.  Outside the framework it reads as the
conditional theorem
\[
  \begin{aligned}
  &\SelBSD\text{-closure}
    \wedge \GammaPD \wedge \GammaSP \wedge \GammaGZ
    \wedge \chiCTp\text{-imports}
    \wedge \AofE\text{-imports} \\
  &\qquad {}\wedge\ \text{Beilinson-imports}
    \Longrightarrow \text{Strong BSD}.
  \end{aligned}
\]
This implication is the paper's result in its plainest form.  The work
below does not erase the hypotheses; it makes their roles explicit and
shows that, once they are supplied, the Strong-BSD scalar identity is
forced by a single closure architecture.

The first component of the architecture is a separation of the deep
arithmetic into two kinds of inputs.  The first kind consists of three
recognition sources.  A recognition source is named arithmetic content
carried by the formed BSD layer: it is supplied as part of the closure
object and used as a hypothesis, not derived from the general framework.
The three sources are \(\GammaPD\), the additive-prime \(p\)-adic
descent/Tamagawa correction; \(\GammaSP\), the signed-Selmer
persistence of the \(\Sha\)-factor through local cover change; and
\(\GammaGZ\), the higher-rank Gross--Zagier-type fixity of the
archimedean regulator readout.  They correspond, respectively, to the
local \(p\)-adic factor, the \(\Sha\)-persistence factor, and the
rank-\(\geq 2\) regulator/fixity factor in the Strong-BSD product.  In
the formal development they are typed carriers of hypotheses, not
free-standing assumptions of the framework; in the paper they are treated
as named recognition content and are introduced before they are used.

The second kind of input consists of imported theorem stacks.  The three
headline imports are the ones used directly in the landing theorem.  The
\(\chiCTp\) stack packages the Cassels--Tate Pfaffian comparison and its
sign convention.  The \(\AofE\) stack packages the Selmer-complex pairing
substrate needed to compare the determinant formulation with the
Cassels--Tate side.  The Beilinson stack packages the rank-\(\geq 2\)
archimedean determinant--\(L\)-value comparison.  Later sections also
record the overconvergent \(\eta\) and \(T_{\mathrm{CASCADE}}\) stacks as
supporting conditional imports, but those are not hidden inside the three
headline imports.  These imports name classical arithmetic geometry as
provenance: Sakamoto, Macias Castillo--Sano, Nekov\'a\v{r}, Flach, the
Cassels--Tate literature, Burns--Flach, Beilinson, and the surrounding
Selmer-complex and determinant technology.  They are not reproved here.
They enter as proof-carrying assumptions, and later sections keep their
individual roles separate.

The second component is the saturated BSD trace shell.  It is the
five-part object
\[
  \SelBSD=(E,\rhoE,J^{!}_{\mathrm{BSD}},
  \mathrm{Vis}^{!}_{\mathrm{BSD}},\mathrm{Audit}^{!}_{\mathrm{BSD}}),
\]
where \(E/\Q\) is the elliptic curve, \(J^{!}_{\mathrm{BSD}}\) records
the duality or involutive datum, \(\mathrm{Vis}^{!}_{\mathrm{BSD}}\)
records the visible scalar readout, and
\(\mathrm{Audit}^{!}_{\mathrm{BSD}}\) records the closure audit.  The
central scalar is the residual
\[
  \rhoE
  =
  \frac{L^{(r)}(E,1)}{r!}
  -
  \frac{\#\Sha(E/\Q)\,\Reg_{\mathrm{NT}}(E)\,\Omega_E
    \prod_v c_v(E)}
       {|E(\Q)_{\tors}|^2}.
\]
Thus \(\rhoE=0\) is not a new BSD-like statement: it is exactly the
standard Strong-BSD scalar identity written as a residual equation.
The point of the shell is to keep the analytic side, the arithmetic
product, the duality datum, the visible readout, and the audit data in
one typed object, so that later arguments can say precisely which input
changes which part of the scalar package.

The recognition predicate \(\PiBSD\) is the shell's structured readout.
It is not defined to be Strong BSD.  It is a quantified predicate that
asks for the three recognition-source readouts in their proper ranges:
the \(p\)-adic descent readout for every additive bad prime and
trivializing local lift in scope, the signed-Selmer persistence readout
for every admissible signed-Selmer prime, and the higher-Gross--Zagier
fixity readout in rank at least two, with the classical low-rank lane
for ranks \(0\) and \(1\).  The readout-level equivalence says that
vanishing of the shell residual is equivalent to the usual Strong-BSD
scalar identity.  This equivalence is a genuine algebraic step: it
rewrites the residual, rather than making \(\PiBSD\) a synonym for the
conjecture.

The headline landing then has a short conceptual chain.  The recognition
sources and imports are supplied.  The higher-rank fixity source forces
the shell residual to vanish, after the \(\kapr\)-normalization identifies
the cascade regulator product with the standard
\(\#\Sha/|E(\Q)_{\tors}|^2\) factor.  The readout equivalence converts
\(\rhoE=0\) into
\[
  \frac{L^{(r)}(E,1)}{r!}
  =
  \frac{\#\Sha(E/\Q)\,\Reg_{\mathrm{NT}}(E)\,\Omega_E
    \prod_v c_v(E)}
       {|E(\Q)_{\tors}|^2}.
\]
The master-theorem applicability audit checks that the three recognition
sources have the correct framework profile, and the composite signature
checks that the three source-governed factors are not disposable: changing
the \(\Sha\)-factor, the regulator factor, or the Tamagawa factor changes
the scalar package.  This is the sense in which the conclusion is a
closure result rather than a restatement of its target.

Two ingredients are imported from \citet{TsiokosBSDApparatus2026}.
First, the \(\kapr\) normalization in that work is the bridge from the cascade leading-term
form to the standard Strong-BSD factor \(\#\Sha/|E(\Q)_{\tors}|^2\).
Second, its adequacy and five-column framing provide the
diagnostic language for separating height, period, finite, \(p\)-adic,
and determinant contributions.  This article uses those results as cited
inputs; it does not refer to internal labels of the cited work.  The AOR
terminology used at the end of the paper is the Audited Operational
Realisability language of \cite{TsiokosAOR2026}, and the recognition-mode
discipline is part of the audited interaction calculus of
\cite{TsiokosFoundationsIII2026}.

The final move is a reading of the completed carrier as a refinement-stable
audited object.  Audited Operational Realisability is a refinement-stable,
audit-closed membership criterion: informally, an object qualifies when
the residuals left by refinement have been accounted for by approved
mechanical records, imported records, or bridged recognition records.  In
this paper the object is \(\SelBSD\).  The resulting statement is the
framework's ``Strong BSD holds in reality'' reading, but it is
non-eliminative: the recognition records are reclassified as bridged
records, not made to disappear.

We close the introduction by fixing the scope.  This is not an
unconditional classical proof of BSD.  It is a conditional closure under
the standing shell-closure record, the three recognition sources, and the
headline imported theorem stacks displayed above.  The rank-\(\geq 2\)
generalized Gross--Zagier identity
\(T_{E12}\) and the additive \(p\leq 3\) Iwasawa main conjecture
\(T_{\mathrm{BAD}}\) remain open audit-result no-gos: the paper records
where they are needed and how the recognition sources carry them, but
does not prove them.  Likewise, the AOR-instance reading is
non-eliminative; it records an audit-closed presentation of the conditional
closure, not the discharge of every open arithmetic obligation.

\section{The imported theorem stacks}
\label{sec:imports}

The closure separates its inputs into two classes.  The present section
records the imported theorem stacks: established arithmetic geometry
that the closure consumes as proof-carrying hypotheses.  The next section
records the recognition sources: arithmetic content carried by the
formed BSD layer.  This distinction matters.  The imported stacks are not
new theorems here, and the closure never derives them
from the framework.  It derives consequences from them, after their
provenance and downstream role have been made explicit.

There are three headline imports in this section.  The first controls the
Cassels--Tate side of the finite \(p\)-primary comparison.  The second
supplies the Selmer-complex pairing substrate \( \AofE \).  The third
supplies the rank-\(\geq 2\) archimedean determinant comparison.  Together
they keep the classical arithmetic geometry used by the landing theorem out
of the framework argument while still making clear exactly which classical
inputs are being used.  The overconvergent \(\eta\) and
\(T_{\mathrm{CASCADE}}\) imports are introduced later as supporting
conditional stacks and are counted separately in App~\ref{app:formalization}.

The symbol \(\chiCTp\) refers here to a Cassels--Tate comparison: a
comparison between a square-root Fitting object on a Selmer-complex
cohomology group and the Pfaffian scalar attached to the Cassels--Tate
pairing.  In ordinary terms, this is the input that lets the \(p\)-primary
finite piece be read with the correct orientation and sign.

\begin{definition}[Imported $\chi_{CT,p}$ theorem stack]\label{def:closure:chi-ct-p-import}
A proof-carrying structure bundling the eight named external inputs
\(T_{E1}\ldots T_{E8}\): Sakamoto's Stark-system theory~\citep{Sakamoto2018};
the determinant/Stark-system and derived-height comparisons of Macias
Castillo--Sano~\citep[Thms.~3.4 and~4.2]{MaciasCastilloSano2026};
Knudsen--Mumford determinant theory~\citep{KnudsenMumford1976};
Nekov\'a\v{r}'s Selmer-complex pairing~\citep[\S10.8.7]{Nekovar2006};
Flach's pairing comparison~\citep{Flach1990}; and the classical
2-Selmer normalization of Fisher--Schaefer--Stoll~\citep[Thm.~6.3]{FisherSchaeferStoll2010}.
The results are stored as proposition fields,
plus the unconditional sign closure
\(\Sigma_{\mathrm{NekCT}}=+1\).  It is taken as a hypothesis by
\Cref{thm:closure:chi-ct-p-comparison}; it is not derived in this
paper.\footnote{Encoded in Lean as a typed structure carrier
\texttt{SixBirdsBSD.\allowbreak Closure.\allowbreak Imports.\allowbreak chiCTpImport}
bundling the imported theorems as hypothesis fields; not a Lean axiom.
See App~\ref{app:formalization}.}
\end{definition}

The downstream role of this stack is narrow.  It is not a general
replacement for the Cassels--Tate literature.  It is the specific
orientation-respecting package needed for the \(\chiCTp\) theorem:
given the usual finite \(\Sha[p^\infty]\), nondegeneracy, and
Selmer-complex hypotheses, it supplies the Pfaffian comparison and the
sign closure \(\Sigma_{\mathrm{NekCT}}=+1\).

The second import packages the Selmer-complex pairing substrate denoted
\(\AofE\).  This notation is local to this article: \(\AofE\) is the
unified Selmer-complex pairing object, not the analytic leading coefficient
written \(A_E\) in \citet{TsiokosBSDApparatus2026}.  Mathematically, it is the
ambient category, Selmer complex, and pairing target in which the
Cassels--Tate pairing is intrinsic and compatible with the
determinant-line formulation.

\begin{definition}[Imported $\AofE$ Selmer-complex pairing]\label{def:closure:a-e-import}
A proof-carrying structure for the \(\AofE\) unified Selmer-complex
substrate (Nekov\'a\v{r} Selmer complex~\citep{Nekovar2006} with intrinsic Cassels--Tate
pairing; Burns--Flach and related Macias Castillo--Sano--Nekov\'a\v{r}
ETNC-adjacent input~\citep{BurnsFlach2001,BlochKato1990,FontainePerrinRiou1994,BurnsMaciasCastilloSeo2024,MaciasCastilloSano2026}
\(\AofE^{\mathrm{Sel}}\)): a typed target category with a
bilinear/categorical pairing field.  It is used only as a
hypothesis.\footnote{Encoded in Lean as a typed structure carrier
\texttt{SixBirdsBSD.\allowbreak Closure.\allowbreak Imports.\allowbreak aEImport}
bundling the imported theorem as a hypothesis field; not a Lean axiom.
See App~\ref{app:formalization}.}
\end{definition}

This import is the bridge between two ways of seeing the same finite
arithmetic information.  On one side is the Selmer-complex description
with its intrinsic Cassels--Tate pairing.  On the other is the
determinant-line language used in the closure landing.  The \(\AofE\)
carrier records that the target category and pairing data have already
been supplied, so later arguments may use the pairing without
reconstructing the Selmer-complex theory.

The third import is archimedean.  In rank at least two, the landing
needs a comparison between the Beilinson determinant side and the
leading \(L\)-value side.  This is where the regulator-to-determinant
identification enters.  The closure treats it as a classical imported
comparison, separate from the higher-rank recognition source that will
encode the remaining Gross--Zagier-type fixity.

\begin{definition}[Imported Beilinson rank-$\geq 2$ comparison]\label{def:closure:beilinson-import}
A proof-carrying structure for the Beilinson determinant/\(L\)-value~\citep{Beilinson1985}
comparison in rank \(r\geq 2\): the archimedean
regulator-to-determinant-line identification used by the rank-\(\geq 2\)
landing.  It is used only as a hypothesis.\footnote{Encoded in Lean as a
typed structure carrier
\texttt{SixBirdsBSD.\allowbreak Closure.\allowbreak Imports.\allowbreak beilinsonImport}
bundling the imported theorem as a hypothesis field; not a Lean axiom.
See App~\ref{app:formalization}.}
\end{definition}

The point of isolating this comparison is to prevent the closure
argument from hiding an archimedean theorem inside the recognition
predicate.  The Beilinson import supplies the determinant/\(L\)-value
identification; the recognition source introduced later supplies the
formed-layer fixity condition that makes the rank-\(\geq 2\) readout
land on the Strong-BSD scalar.

\begin{nonclaim}\label{nonclaim:closure:imports-not-derived}
The \(\chiCTp\) comparison remains conditional on its imported stack
\(T_{E1}\ldots T_{E8}\); the import structures are sourced assumptions,
not results of this paper.  The closure derives consequences
\emph{from} them and never proves them.  They are carried as typed
hypothesis structures and introduce no project-local axioms.
\end{nonclaim}

% ===========================================================================

\section{The recognition hypotheses}
\label{sec:recognition-sources}
% Each carrier is a typed structure in Lean whose arithmetic content is a
% hypothesis field (lean_coverage = recognition_source).

In the Six Birds framework a \emph{recognition source} is arithmetic
content that a construction rests on but does not derive: it is stated
precisely, attached to the object, and used as a hypothesis.  In this
paper there are three, one for each of three factors of the BSD formula
that the formal construction does not supply by itself.
\begin{itemize}
\item $\GammaPD$ concerns the local Tamagawa factor at an odd prime of
  additive reduction;
\item $\GammaSP$ concerns the $\Sha$ factor, through a signed local
  comparison;
\item $\GammaGZ$ concerns the leading term in analytic rank at least $2$.
\end{itemize}
\Cref{tab:recognition-sources} summarizes them.  We say at once what will
matter for the main theorem: its scalar conclusion uses only $\GammaGZ$ in
rank at least $2$, and the classical identity in rank at most $1$
(\Cref{sec:shell-readout}).  The two local hypotheses are part of the
recognition predicate, with their scopes recorded below, but the scalar
argument does not use them.
\begin{table}[tbp]
\centering
\footnotesize
\caption{The three recognition hypotheses.  The last column records
whether the scalar argument of \Cref{thm:closure:strong-bsd-conditional}
uses the hypothesis.}
\label{tab:recognition-sources}
\begin{tabularx}{\textwidth}{@{}>{\raggedright\arraybackslash}p{0.17\textwidth}
>{\raggedright\arraybackslash}p{0.25\textwidth}
>{\raggedright\arraybackslash}X
>{\raggedright\arraybackslash}p{0.12\textwidth}@{}}
\toprule
Hypothesis & Scope & Content & Used for scalar BSD \\
\midrule
\(\GammaPD\) & odd prime \(p\), additive type II or III, \(c_p\in\{1,2\}\),
\(\Sha[p^\infty]=0\), tame \(L/\Qp\) giving good reduction & descent congruence
modulo \(\Zp^\times\) with factor \(c_p^{-1}\) and a unit; a local factor
matched with \(\Tam(E)\) & no \\
\addlinespace
\(\GammaSP\) & prime \(p\) of additive reduction in the signed scope, tame \(L/\Qp\) giving good reduction & signed corestriction determinant \(\equiv c_p\)
modulo \(\Zp^\times\); a factor matched with \(\#\Sha(E/\Q)\) & no \\
\addlinespace
\(\GammaGZ\) & no CM, analytic rank \(r\ge2\) & \(\psi_-(E,r)=0\), with
\(\kapr=\#\Sha/\#E(\Q)_{\tors}^{\,2}\) and matched factors; equivalent to
Strong BSD for \(E\) when \(\Sha\) is finite & yes, in rank \(\ge2\) \\
\bottomrule
\end{tabularx}
\end{table}


\begin{definition}[$\GammaPD$: $p$-adic descent hypothesis]\label{def:closure:gamma-padic-descent}
Let $E/\Q$ and let $p$ be an odd prime at which $E$ has additive
reduction of Kodaira type II or III, with $c_p(E)\in\{1,2\}$, with
$\Sha(E/\Q)[p^\infty]=0$, and let $L/\Qp$ be a finite tamely ramified
extension over which $E$ acquires good reduction (equivalently, over
which inertia acts trivially on the $\ell$-adic Tate module for every
prime $\ell\neq p$; reduction of type II or III is potentially good).  The hypothesis
$\GammaPD$ at $(E,p,L)$ supplies a descent map
\[
  \mathrm{descent}_p^{\mathrm{OC}}:
  LC_{\mathrm{triv}}(L_p^{\mathrm{OC}}(E_L))
  \longrightarrow
  LC_{\mathrm{triv}}(L_p^{\mathrm{OC}}(E))
\]
between trivialized leading coefficients of ordinary/anticyclotomic local
data over $L$ and over $\Qp$, a unit $\varepsilon_p(E,L)\in\Zp^\times$, and
the congruence
\[
  LC_{\mathrm{triv}}(L_p^{\mathrm{OC}}(E)) \equiv
  c_p(E)^{-1}\,\varepsilon_p(E,L)\,
  \mathrm{descent}_p^{\mathrm{OC}}
    \bigl(LC_{\mathrm{triv}}(L_p^{\mathrm{OC}}(E_L))\bigr)
  \pmod{\Zp^\times}.
\]
It also carries a factor $\mathrm{tam}_p$, which the recognition predicate
of \Cref{sec:shell-readout} matches with the Tamagawa product of $E$.
\end{definition}

Two features of this hypothesis should be noted.  First, since $p$ is odd
and $c_p(E)\in\{1,2\}$, the number $c_p(E)$ is a $p$-adic unit (in
particular $p\nmid c_p(E)$).  The congruence modulo $\Zp^\times$ therefore
holds or fails independently of whether $c_p(E)$ is $1$ or $2$: it carries
no information about $c_p(E)$.  For example, the curves $121a1$ and
$121b1$ have additive reduction at $11$ of types II and III, with $c_{11}$
equal to $1$ and $2$ respectively, and these two values have the same
class modulo $\Z_{11}^\times$.  The value of $c_p(E)$ is recorded in the
hypothesis separately.  Second, the matching of a local factor with the
global Tamagawa product is itself a comparison.  On the curve
$y^2=x^3+385x+1225$, the two places $5$ and $7$ of type III each have
$c_p=2$, and the global Tamagawa product is $4$; a single local factor is
not the global one.

\begin{definition}[$\GammaSP$: signed $\Sha$-persistence hypothesis]\label{def:closure:gamma-sha-persistence}
Let $E/\Q$, let $p$ be a prime of additive reduction at which the signed
Selmer theory used here applies, and let $L/\Qp$ be a finite tamely
ramified extension over which $E$ acquires good reduction.  Let
$\mathrm{sp}_p(E,L)$ be the signed local pairing on the cover side,
$\cores_{L/\Qp}$ the corestriction from the cover-side target to a
base-side target, and
\[
  C_{\mathrm{loc}}^{\mathrm{III}}(E,p)
  =
  \det\bigl(\cores_{L/\Qp}\circ \mathrm{sp}_p(E,L)\bigr)
\]
the determinant of the composite.  This construction, its
determinant-line convention and its applicability at $(E,p,L)$ are part
of the hypothesis; the signed theory of Kobayashi and
Pollack~\citep{Kobayashi2003,Pollack2003}, developed for good
supersingular reduction, is background for it, not a source of it.  The
hypothesis $\GammaSP$ at $(E,p,L)$ asserts
\[
  C_{\mathrm{loc}}^{\mathrm{III}}(E,p)
  \equiv c_p(E)\pmod{\Zp^\times},
\]
and carries a factor $\mathrm{sha}_p$, which the recognition predicate
matches with $\#\Sha(E/\Q)$.
\end{definition}

The same remark applies: when $p\nmid c_p(E)$, the congruence does not
see the value of $c_p(E)$.  The determinant of a composite of a map into
the cover-side target with corestriction to the base side requires a
determinant-line convention, which is part of the hypothesis, as is the
matching of $\mathrm{sha}_p$ with the global order of $\Sha$.

\begin{definition}[$\GammaGZ$: leading-term hypothesis in rank at least two]\label{def:closure:gamma-higher-gz-fixity}
Let $E/\Q$ be without complex multiplication, of analytic rank $r\ge2$.
Put
\[
  \psi_-(E,r)
  =
  \frac{L^{(r)}(E,1)}{r!}
  -
  \kapr(E)\,\Reg_{\mathrm{NT}}(E)\,\Omega_E\,\Tam(E).
\]
The hypothesis $\GammaGZ$ at $E$ asserts $\psi_-(E,r)=0$.  In the
recognition predicate it is used together with the normalization
$\kapr(E)=\#\Sha(E/\Q)/\#E(\Q)_{\tors}^{\,2}$ and with the matching of
its regulator, period and Tamagawa factors to those of $E$.
\end{definition}

The name refers to the kind of theorem that would be expected to supply
this hypothesis, a generalized Gross--Zagier formula in rank at least
$2$.  No theorem supplying this normalized higher-rank identity is
established among the inputs used here, and we retain it as a hypothesis
(see also \Cref{rmk:closure:t-e12-gap}).  It is
important to be clear about the strength of $\GammaGZ$.  Assume
$\Sha(E/\Q)$ finite, and let $\Omega_E$ and $\Tam(E)=\prod_vc_v(E)$ be as
in the BSD formula.  By the normalization of $\kapr$
\citep[Section~13]{TsiokosBSDApparatus2026}, the arithmetic term of
$\psi_-(E,r)$ is then
\[
  \kapr(E)\,\Reg_{\mathrm{NT}}(E)\,\Omega_E\,\Tam(E)
  =
  \frac{\#\Sha(E/\Q)\,\Reg_{\mathrm{NT}}(E)\,\Omega_E\,\Tam(E)}
       {\#E(\Q)_{\tors}^{\,2}},
\]
so $\GammaGZ$, read with this normalization, is the Strong BSD identity for
$E$.  The hypothesis is not weaker than the conclusion it is used for.

\begin{nonclaim}\label{nonclaim:closure:recognition-sources-not-derived}
None of the three recognition hypotheses is derived in this paper, from
the Six Birds framework or otherwise.  Each is a hypothesis about a given
curve, recorded in the formalization as a structure whose arithmetic
content is a field, not as an axiom.  We also do not claim that any of
them is necessary for the scalar conclusion: in rank at least $2$,
$\GammaGZ$ with its normalization already implies it.
\end{nonclaim}

\section{The saturated trace shell and its readout}
\label{sec:shell-readout}

The preceding section isolated the three arithmetic recognition sources.
We now assemble the object on which they act.  The saturated BSD trace
shell is a typed carrier for the Strong-BSD scalar comparison: it keeps
the analytic leading coefficient, the arithmetic product, and the audit
data in one place, so that the closure can record which input governs
which scalar factor.  Here ``trace shell'' means a formed mathematical
object together with its visible scalar readout and the bookkeeping
needed to prevent the readout from being a tautological restatement of
the desired identity.

\begin{definition}[Saturated BSD trace shell]\label{def:closure:sel-bsd-shell}
\(\SelBSD\) is the five-tuple
\[
  \SelBSD
  =
  \bigl(E,\ \rhoE,\ J^{!}_{\mathrm{BSD}},\
    \mathrm{Vis}^{!}_{\mathrm{BSD}},\
    \mathrm{Audit}^{!}_{\mathrm{BSD}}\bigr).
\]
Here \(E/\Q\) is the elliptic curve and \(\rhoE\) is the scalar
residual
\[
  \rhoE
  =
  \frac{L^{(r)}(E,1)}{r!}
  -
  \frac{\#\Sha(E/\Q)\,\Reg_{\mathrm{NT}}(E)\,\Omega_E\,\Tam(E)}
       {|E(\Q)_{\tors}|^2}.
\]
Thus \(\rhoE=0\) is exactly the Strong-BSD scalar identity, with
\(\Tam(E)\) read as the Tamagawa product.  The datum
\(J^{!}_{\mathrm{BSD}}\) records the involutive or duality structure
that compares the analytic and arithmetic sides.  The datum
\(\mathrm{Vis}^{!}_{\mathrm{BSD}}\) records the visible scalar readout
of the shell.  The datum \(\mathrm{Audit}^{!}_{\mathrm{BSD}}\) records
the admissibility, anti-tautology, and lawfulness information attached
to the formed object.  The standing closure assumption used here is the
standard assumption that the formed object carries this fixed content
itself: admissibility, anti-tautology, and lawfulness data are part of
the shell rather than consequences inferred after the fact.
\end{definition}

\Cref{fig:sel-bsd-shell-schematic,tab:sel-bsd-shell-fields} summarize the
shell fields and the residual readout used below.
\begin{figure}[tbp]
\centering
\begin{tikzpicture}[
  font=\small,
  box/.style={draw, rounded corners=2pt, align=center, minimum width=2.5cm,
    minimum height=0.85cm, fill=black!3},
  shell/.style={draw, rounded corners=3pt, align=center, minimum width=4.2cm,
    minimum height=1.1cm, fill=black!8},
  arrow/.style={-{Latex[length=2mm]}, thick}
]
\node[shell] (shell) at (0,0) {\(\SelBSD=(E,\rhoE,J^{!},\mathrm{Vis}^{!},\mathrm{Audit}^{!})\)};
\node[box] (analytic) at (-4.6,1.7) {analytic\\leading term};
\node[box] (arith) at (0,1.7) {arithmetic\\product};
\node[box] (audit) at (4.6,1.7) {audit\\records};
\node[box] (readout) at (0,-1.8) {\(\rhoE=0\)\\readout};

\draw[arrow] (analytic) -- (shell);
\draw[arrow] (arith) -- (shell);
\draw[arrow] (audit) -- (shell);
\draw[arrow] (shell) -- (readout);
\node[align=center, font=\footnotesize] at (0,-2.75)
  {the readout is equivalent to the Strong-BSD scalar identity};
\end{tikzpicture}
\caption{The saturated BSD trace shell packages the residual, readout data, and
audit records used by the conditional closure.}
\label{fig:sel-bsd-shell-schematic}
\end{figure}

\begin{table}[tbp]
\centering
\small
\caption{Fields of the saturated BSD trace shell.}
\label{tab:sel-bsd-shell-fields}
\begin{tabularx}{\textwidth}{@{}>{\raggedright\arraybackslash}p{0.18\textwidth}
>{\raggedright\arraybackslash}p{0.34\textwidth}
>{\raggedright\arraybackslash}X@{}}
\toprule
Field & Role & Paper-prose name \\
\midrule
\(E\) & The elliptic curve over \(\Q\). & base curve \\
\(\rhoE\) & Difference between the analytic leading coefficient and the
Strong-BSD arithmetic product. & scalar residual \\
\(J^{!}_{\mathrm{BSD}}\) & Involutive or duality datum comparing analytic and
arithmetic sides. & duality datum \\
\(\mathrm{Vis}^{!}_{\mathrm{BSD}}\) & Visible scalar readout data attached
to the shell. & visible readout \\
\(\mathrm{Audit}^{!}_{\mathrm{BSD}}\) & Admissibility, anti-tautology, and
lawfulness records. & audit record \\
\addlinespace
\(\rhoE=0\) & The residual vanishing equation, equivalent to the Strong-BSD
scalar identity at readout level. & landing readout \\
\bottomrule
\end{tabularx}
\end{table}


The predicate placed on the shell is not the Strong-BSD identity by
definition.  It is a structured readout predicate: it asks for the three
recognition-source statements, quantified over the local and rank data
where they apply.  This distinction is essential.  The next theorem can
then identify the readout with the Strong-BSD scalar identity; it is not
merely unfolding a definition.

\begin{definition}[$\Pi_{\mathrm{BSD}}$ predicate]\label{def:closure:pi-bsd}
\(\PiBSD(E)\) is the conjunction of the three recognition-source
readouts on \(\SelBSD\), with the relevant quantifiers retained.  First,
for every additive bad prime \(p\) in scope and every tame
inertia-trivializing local lift \(L/\Qp\) in scope, the
\(\GammaPD\) readout supplies the \(p\)-adic descent/Tamagawa
correction.  Second, for every signed-Selmer-admissible prime \(p\),
the \(\GammaSP\) readout supplies the signed-Selmer local determinant
matching.  Third, for analytic rank \(r\geq 2\), the \(\GammaGZ\)
readout supplies the higher-rank fixity assertion
\(\psi_-(E,r)=0\); for \(r\in\{0,1\}\), the predicate uses the
classical low-rank lane.  In particular, \(\PiBSD\) is a structured
quantified predicate over recognition readouts, not the definition
\(\PiBSD:=\mathrm{StrongBSD}\).
\end{definition}

\begin{theorem}[Readout-level equivalence]\label{thm:closure:pi-bsd-iff-strong-bsd}
The scalar readout of \(\PiBSD\) on \(\SelBSD\) is equivalent to the
standard Strong-BSD scalar identity
\[
  \frac{L^{(r)}(E,1)}{r!}
  =
  \frac{\#\Sha(E/\Q)\,\Reg_{\mathrm{NT}}(E)\,\Omega_E\,
    \prod_v c_v(E)}
       {|E(\Q)_{\tors}|^2}.
\]
This is a readout-level equivalence: it identifies the visible scalar
readout with the Strong-BSD scalar identity, but it does not define
\(\PiBSD\) to be Strong BSD and it does not manufacture the recognition
sources from the scalar identity alone.\footnote{Verified in Lean as
\texttt{SixBirdsBSD.\allowbreak Closure.\allowbreak SelShell.\allowbreak piBSDIffStrongBSD};
see App~\ref{app:formalization}.}
\end{theorem}

\begin{proof}
By definition of the shell, the visible scalar residual is
\[
  \rhoE
  =
  \frac{L^{(r)}(E,1)}{r!}
  -
  \frac{\#\Sha(E/\Q)\,\Reg_{\mathrm{NT}}(E)\,\Omega_E\,\Tam(E)}
       {|E(\Q)_{\tors}|^2}.
\]
The readout predicate \(\PiBSD\) supplies the three pieces entering this
residual.  The higher-rank source \(\GammaGZ\) gives
\(\psi_-(E,r)=0\).  After applying the \(\kapr\) normalization of
\citet{TsiokosBSDApparatus2026},
\[
  \kapr(E)=\frac{\#\Sha(E/\Q)}{|E(\Q)_{\tors}|^2},
  \qquad
  \kapr(E)\,\Reg_{\mathrm{NT}}(E)\,\Omega_E\,\Tam(E)
  =
  \frac{\#\Sha(E/\Q)\,\Reg_{\mathrm{NT}}(E)\,\Omega_E
    \prod_v c_v(E)}
       {|E(\Q)_{\tors}|^2}.
\]
Thus the equation \(\psi_-(E,r)=0\) is exactly the assertion
\(\rhoE=0\) in rank at least two.  The source \(\GammaPD\) supplies the
local \(p\)-adic/Tamagawa correction entering the Tamagawa product, and
\(\GammaSP\) supplies the signed-Selmer persistence of the \(\Sha\)
factor; these readouts fix the remaining scalar factors in the same
residual.  Therefore the forward scalar readout of \(\PiBSD\) is
\(\rhoE=0\).

Finally, the equation \(\rhoE=0\) is algebraically the displayed
Strong-BSD scalar identity: moving the arithmetic product to the other
side gives the identity, and subtracting the right-hand side from the
left gives \(\rhoE=0\).  Conversely, if the displayed scalar identity is
already known as a scalar readout, then the same subtraction gives
\(\rhoE=0\).  This reverse implication is only an equivalence of
scalar readouts; it does not construct any of the three
\(\Gamma_{\mathrm{BSD}}^{*}\) predicates.
\end{proof}

\begin{theorem}[Master-theorem applicability]\label{thm:closure:master-theorem-applicability}
The closure master theorem applies to \(\SelBSD\): for the audited
interaction calculus \cite{TsiokosFoundationsIII2026}, correctly
instantiated recognition sources that exhaust the residual factors, and
that pass the no-smuggling audit, force the residual readout
\(\rhoE=0\).\footnote{Verified in Lean as
\texttt{SixBirdsBSD.\allowbreak Closure.\allowbreak SelShell.\allowbreak masterTheoremApplicability};
see App~\ref{app:formalization}.}
\end{theorem}

\begin{proof}
We verify the two applicability requirements for the shell.  First, the
recognition sources are instantiated on exactly the scalar factors of
\(\rhoE\).  The source \(\GammaPD\) governs the local
\(p\)-adic/Tamagawa correction: it supplies the descent congruence and
the Tamagawa factor \(c_p(E)\) up to a \(p\)-adic unit.  The source
\(\GammaSP\) governs the persistence of the \(\Sha\)-factor under the
signed-Selmer cover change: its determinant comparison is the local
component entering the same arithmetic product.  The source
\(\GammaGZ\) governs the rank-\(r\) regulator/fixity factor: it states
that the analytic leading coefficient equals the arithmetic regulator
product after the \(\kapr\) normalization.  These three pieces account
for the local \(p\)-adic/Tamagawa factor, the \(\Sha\)-persistence
factor, and the regulator/fixity factor respectively.  No scalar factor
of \(\rhoE\) remains outside this list, and no factor is assigned to two
different sources.

Second, the no-smuggling condition holds.  The imported theorem stacks
enter only as hypotheses, and the recognition sources enter only as
carried readouts.  The closure does not silently turn an open
arithmetic obligation into a derivation target; it records the
obligation as supplied content and then derives consequences from it.
Thus the recognition-source instantiation is exhaustive for the
residual and the audit contains no hidden derivation of out-of-scope
arithmetic.  The closure master theorem is therefore applicable to
\(\SelBSD\), and its residual conclusion is \(\rhoE=0\).
\end{proof}

\begin{theorem}[Composite signature with anti-tautology hardening]\label{thm:closure:composite-signature}
The composite recognition predicate attached to \(\SelBSD\) is
computable and falsifiable, carries a framework trace signature as a
fingerprint of the composite operator, and is anti-tautological: it
genuinely consumes all three recognition records.  In particular, the
three components are independent but not disposable.  Ablating
\(\GammaPD\), \(\GammaSP\), or \(\GammaGZ\) breaks the Strong-BSD
scalar factor package, so the composite does not collapse to a
tautological restatement of Strong BSD.\footnote{\raggedright Tracked in the
formalization harness as a conditional schema
\texttt{SixBirdsBSD.\allowbreak Closure.\allowbreak SelShell.\allowbreak compositeSignature}
-- a checked projection over an explicit proof-carrying record
(conditional on the imports and recognition sources), not an independent
Lean derivation. See App~\ref{app:formalization}.}
\end{theorem}

\begin{proof}
The composite predicate is a finite conjunction of the three readouts
introduced above, together with the shell data on which those readouts
are evaluated.  It is therefore computable in the framework sense: once
the shell and the recognition records are supplied, the predicate has a
definite list of components to inspect.  It is also falsifiable: failure
of any one component is a failure of the composite readout.

The substantive point is non-disposability.  Removing \(\GammaPD\)
leaves the local \(p\)-adic/Tamagawa factor undetermined; the shell no
longer knows that the descent comparison repairs the Tamagawa
coarsening by the required factor.  Removing \(\GammaSP\) leaves the
\(\Sha\)-persistence factor undetermined; the signed-Selmer determinant
comparison no longer supplies the component entering the arithmetic
product.  Removing \(\GammaGZ\) leaves the rank-\(r\)
regulator/fixity factor undetermined; there is then no assertion that
the analytic leading coefficient agrees with the normalized regulator
product.  Each recognition source governs a different scalar factor of
the residual, and each factor is required for the Strong-BSD scalar
package.  Hence the composite predicate depends genuinely on the
formed BSD layer's data and cannot be replaced by the bare sentence
``Strong BSD holds.''
\end{proof}

% ===========================================================================

\section{The Cassels--Tate comparison}
\label{sec:chi-ct-p}

The first supporting comparison used by the closure is the
Cassels--Tate Pfaffian comparison~\citep{Cassels1962,PoonenStoll1999}.  The import stack introduced in
Definition~\ref{def:closure:chi-ct-p-import} packages the relevant external
results as hypotheses: the Selmer-complex inputs of Sakamoto~\citep{Sakamoto2018}
and Macias Castillo--Sano~\citep{MaciasCastilloSano2026}, the determinant
and orientation inputs of Knudsen--Mumford and Nekov\'a\v{r}, and Flach's
pairing comparison.  Fisher--Schaefer--Stoll~\citep{FisherSchaeferStoll2010}
fix the classical 2-Selmer pairing normalization; they are not used as a
Fitting-ideal theorem.  This section states the
comparison in the form consumed by the closure.  It is conditional on
that stack; no part of the argument below reproves the imported
arithmetic geometry.

\begin{theorem}[$\chi_{CT,p}$ comparison (v5)]\label{thm:closure:chi-ct-p-comparison}
Let \(E/\Q\) be an elliptic curve such that \(\Sha(E)[p^\infty]\) is
finite, the Cassels--Tate pairing is nondegenerate, the standard
Sakamoto and Macias Castillo--Sano Selmer-complex hypotheses hold, and
\(p\nmid \Tam(E/\Q)\).  Conditional on the imported \(\chiCTp\)
structure, there is a canonical orientation-respecting Pfaffian
comparison map
\[
  \Pf_{\mathrm{Nek},p}^{\mathrm{or},\sqrt{\ }}:
  \sqrt{\Fitt}_{\mathrm{CT}}
    \bigl(H^1(C_p(E))/\mathrm{div},\,U_{2,2}\bigr)
  \longrightarrow
  \mathrm{vctp}_p
\]
from the square-root Cassels--Tate Fitting ideal on the quotient by the
divisible subgroup, paired by the Nekov\'a\v{r} cup-product
\(U_{2,2}\) identified with the Flach pairing, to the
Cassels--Tate Pfaffian scalar line \(\mathrm{vctp}_p\).  This map sends
the square-root Stark generator~\citep{MazurRubin2004} \(\sqrt{F}_E^p\) to the
Cassels--Tate height \(h_p^{\mathrm{CT}}(E)\), and the Nekov\'a\v{r}
sign convention satisfies \(\Sigma_{\mathrm{NekCT}}=+1\).\footnote{Tracked in the formalization harness as a conditional schema
\texttt{SixBirdsBSD.\allowbreak Closure.\allowbreak ChiCTp.\allowbreak chiCTpComparison}
-- a checked projection over an explicit proof-carrying record
(conditional on the imported \(\chiCTp\) stack), not an independent Lean
derivation. See App~\ref{app:formalization}.}
\end{theorem}

\begin{proof}
Assume the imported \(\chiCTp\) structure.  Its named components supply
the ingredients in the statement.  The Selmer-complex hypotheses and
the Cassels--Tate nondegeneracy input place the quotient
\(H^1(C_p(E))/\mathrm{div}\) in the self-dual setting in which a
Pfaffian comparison is defined.  The determinant-line and orientation
inputs identify the square-root Fitting ideal
\(\sqrt{\Fitt}_{\mathrm{CT}}(H^1(C_p(E))/\mathrm{div},U_{2,2})\) as the
source of the comparison map.  Here the square-root Fitting ideal is
the square-root determinant object attached to the Cassels--Tate
pairing; the notation \(U_{2,2}\) records the Nekov\'a\v{r}
cup-product, and the imported Flach comparison identifies it with the
classical Cassels--Tate pairing used on the target side.

The Pfaffian input in the stack then gives the canonical
orientation-respecting map
\[
  \Pf_{\mathrm{Nek},p}^{\mathrm{or},\sqrt{\ }}:
  \sqrt{\Fitt}_{\mathrm{CT}}
    \bigl(H^1(C_p(E))/\mathrm{div},U_{2,2}\bigr)
  \longrightarrow
  \mathrm{vctp}_p .
\]
The square-root Stark-generator identification supplies the image
formula: the distinguished generator \(\sqrt{F}_E^p\) maps to the
Cassels--Tate height \(h_p^{\mathrm{CT}}(E)\).  Finally, the sign
closure component of the same imported record fixes the orientation
sign as \(\Sigma_{\mathrm{NekCT}}=+1\).  Combining these imported
components proves exactly the displayed comparison.  The proof is
therefore an assembly under the stated hypotheses; it is precisely as
strong as the imported theorem stack and does not rederive the
Sakamoto, Macias Castillo--Sano, Nekov\'a\v{r}, Flach, or classical
Cassels--Tate comparison inputs.
\end{proof}

% ===========================================================================

\section{The overconvergent \texorpdfstring{$\eta$}{eta}-formula}
\label{sec:eta-formula}

This section records a proposed $p$-adic identity at additive odd primes,
the overconvergent $\eta$-formula.  It plays no role in the main theorem.
We state it only as a conditional input, and we explain why the published
results one might hope to derive it from do not apply in the form needed.

\begin{definition}[$\eta$-certificate]\label{def:closure:eta-published-imports}
Let $E/\Q$ and let $p$ be an odd prime at which $E$ has additive reduction.
An \emph{$\eta$-certificate} at $(E,p)$ consists of: a nonzero $\beta$
(the proposed $U_p$-eigenvalue of an additive $p^2$-stabilization) with
$\beta^2=-p$; a leading term $T_{\mathrm{lead}}(E,p)$ of an overconvergent
expansion, with its order $\mathrm{ord}_T(E,p)\in\Z$; a value
$\etap(E)$; nonzero scalars standing for $p$ and $\#E(\Q)_{\tors}$; and
the identity
\[
  \etap(E)
  =
  \bigl(p\,\#E(\Q)_{\tors}\bigr)^{\mathrm{ord}_T(E,p)-1}
  \,\beta^{-2}\,
  T_{\mathrm{lead}}(E,p).
\]
The certificate also lists, as propositions, the background results it
is modelled on: Bella\"iche's critical-slope theory and moment
lemma~\citep{Bellaiche2012}, overconvergent modular
symbols~\citep{PollackStevens2011,ColemanMazur1998}, critical $p$-adic
$L$-functions~\citep{BenoisBuyukboduk2024}, the Eisenstein ideal at
prime-square level~\citep{LangWake2025}, component groups of N\'eron
models~\citep{Lorenzini1993,Ling1997}, and a nonvanishing and primitivity
statement for a cuspidal evaluation, which is a conjecture.
\end{definition}

\begin{theorem}[Conditional $\eta$-formula]\label{thm:closure:oc-eta-formula}
If $(E,p)$ carries an $\eta$-certificate, then
\[
  \etap(E)
  =
  \bigl(p\,\#E(\Q)_{\tors}\bigr)^{\mathrm{ord}_T(E,p)-1}
  \,\beta^{-2}\,
  T_{\mathrm{lead}}(E,p).
\]%
\footnote{The formalization records this as a conditional schema: the
identity is a field of the certificate and the conclusion is its
projection.}
\end{theorem}

\begin{proof}
The identity is part of the certificate.
\end{proof}

We do not claim that the listed background results, together with the
conjecture, imply the identity.  Two elementary obstructions show that at
least some of them do not apply as stated.

First, the relation $\beta^2=-p$ fixes the slope of $\beta$.  With the
valuation normalized by $v_p(p)=1$, it gives $2v_p(\beta)=1$, so
$v_p(\beta)=1/2$.  In weight $2$ the critical slope is $1$; this is the
normalization of Benois--B\"uy\"ukboduk~\citep{BenoisBuyukboduk2024}, for
whom the critical slope in weight $k$ is $k-1$.  So $\beta$ is not of
critical slope, and results about critical-slope eigenvalues do not apply
to it without a further comparison.

Second, the theorem of Lang--Wake~\citep[Theorem~1.1]{LangWake2025}
concerns a level prime $N$ and a residual prime $\ell$, both at least $5$,
with $\ell\mid N+1$; these numerical conditions are necessary, though not
sufficient, for an application.  They exclude the three curves at which
the formula has been proposed, whichever residual prime is chosen.  For
$36a1$ at $p=3$, the level prime would be $3<5$ (and the full level $36$
is not prime).  For $49a1$ at $p=7$, $N+1=8$ has no prime divisor at least
$5$.  For $121b1$ at $p=11$, $N+1=12$ has no prime divisor at least $5$.

Numerical values of $\etap$ at these three curves have been recorded,
but without complete input data or a reproducible computation of the
right-hand side.  We do not regard them as evidence for the formula, at
those curves or in general.%
\footnote{Lean proves the slope computation (from the relation
$\beta^2=-p$, the multiplicativity of the valuation at $\beta^2$ and
$v_p(-p)=1$, it derives $v_p(\beta)=1/2\neq1$) and the exclusion of the
level primes $3$, $7$ and $11$ from the numerical conditions of
Lang--Wake for every residual prime.}

\section{The \texorpdfstring{$p$}{p}-part in rank at most one}
\label{sec:t-cascade}

In analytic rank $0$ and $1$, the recognition predicate contains the
scalar Strong BSD identity itself as a hypothesis
(\Cref{def:closure:pi-bsd}).  This section records, for context, the
form in which known results are expected to bear on that hypothesis: they
give the $p$-part of the identity for suitable primes $p$.  As with the
other inputs, we state precisely what is assumed; the main theorem does not
use this section.

\begin{definition}[Rank-at-most-one input]\label{def:closure:t-cascade-import}
The rank-at-most-one input is a list of forty propositions, each with an
asserted proof, standing for published results in two branches: for curves
with complex multiplication, Rubin's Iwasawa theory over imaginary
quadratic fields~\citep{Rubin1991}, the main conjecture at supersingular
primes of Pollack--Rubin~\citep[Theorem~7.3]{PollackRubin2004}, and work of
Burungale--Castella--Skinner--Tian~\citep{BurungaleCastellaSkinnerTian2022}
at good ordinary primes, together with the archimedean, Heegner-point and
uniformization ingredients needed to compare $p$-adic and classical
quantities; for curves without complex multiplication, the work of
Skinner--Urban~\citep{SkinnerUrban2014}, Wei Zhang~\citep{WeiZhang2014},
Castella et al.~\citep{CHKLL2025} and Yan--Zhu~\citep{YanZhu2026}.  Both
branches are part of the input simultaneously.
\end{definition}

\begin{theorem}[Conditional $p$-part identity in rank at most one]\label{thm:closure:t-cascade-rank-le-1}
Let $E/\Q$ have analytic rank $r\in\{0,1\}$ and let $p$ be a prime.
Assume the hypotheses of one of the two branches:
\begin{itemize}
\item[(CM)] $E$ has complex multiplication by $K$; either $p$ is good
  ordinary and split in $K$, or $p\ge5$ is good supersingular and inert in
  $K$; $\Sha(E/\Q)[p^\infty]$ is finite, the Cassels--Tate pairing is
  nondegenerate, the Selmer-complex hypotheses of Sakamoto and Macias
  Castillo--Sano hold, $p\nmid\Tam(E)$, the standard ramification and
  Galois-image hypotheses hold, and a Heegner point is given when $r=1$;
\item[(non-CM)] $E$ has no complex multiplication, the residual
  representation $\bar\rho_{E,p}$ is surjective~\citep{Serre1972}, $p$ is
  good ordinary with $p\nmid N$ and $p\nmid\Tam(E)$; the local,
  residual-irreducibility, ramification and level hypotheses of
  Skinner--Urban hold; $\Sha(E/\Q)[p^\infty]$ is finite, the Cassels--Tate
  pairing is nondegenerate, the Selmer-complex hypotheses hold, and, when
  $r=1$, the Gross--Zagier--Kolyvagin hypotheses hold for a Heegner field.
\end{itemize}
Assume the rank-at-most-one input, and assume further that, under these
branch hypotheses, the results it lists yield the valuation identity
below.  Then
\[
  v_p\!\left(\frac{L^{(r)}(E,1)/r!}{\Omega_E\,\Reg_{\mathrm{NT}}(E)}\right)
  =
  v_p\!\left(\frac{\#\Sha(E/\Q)\,\Tam(E)}{\#E(\Q)_{\tors}^{\,2}}\right).
\]%
\footnote{The formalization records this as a conditional schema.  Lean
proves that the comparison hypothesis, given the forty asserted
propositions and a curve and prime, is equivalent to the bare implication
from the branch hypotheses to the valuation identity; the theorem is
therefore exactly as strong as that implication.}
\end{theorem}

\begin{proof}
Apply the assumed implication to the branch hypotheses.
\end{proof}

The comparison hypothesis in the theorem is where the arithmetic is: in
the non-CM branch, for instance, one divisibility comes from the
Skinner--Urban main conjecture and the other from Euler systems and the
Gross--Zagier--Kolyvagin method, and making the two match up to $p$-adic
units requires the finiteness, nondegeneracy and nondivisibility
hypotheses.  We have not verified that the listed results yield the
identity under exactly the stated hypotheses, and the input is not
organized curve by curve.  The theorem should be read as a statement of
the expected form of the known $p$-part results, not as a new proof of
them, and it does not supply the full identity required by the
recognition predicate in rank at most one.

\section{Inputs not supplied by known theorems}
\label{sec:obstructions}

Two of the inputs above correspond to well-known open problems.  We record
them as remarks rather than theorems: they describe what, to our
knowledge, is not currently available, and we have not carried out a
systematic search of the literature to support anything stronger.

\begin{remark}[Leading terms in rank at least two]\label{rmk:closure:t-e12-gap}
For $E/\Q$ without complex multiplication and of analytic rank $r\ge2$,
we know of no theorem giving an identity of the form
\[
  \frac{L^{(r)}(E,1)}{r!}
  =
  \kapr(E)\,\det\bigl(\langle Z_i,Z_j\rangle_{\mathrm{NT}}\bigr)\,
  \Omega_E\,\Tam(E)
\]
with explicit classes $Z_1,\dots,Z_r$ forming a $\Z$-basis of
$E(\Q)/E(\Q)_{\tors}$, so that the determinant is $\Reg_{\mathrm{NT}}(E)$
(for points spanning a subgroup of index $m$ it would be
$m^2\Reg_{\mathrm{NT}}(E)$), and explicit normalizing factors, that is, a generalized Gross--Zagier formula in rank at least $2$.  The
hypothesis $\GammaGZ$ is exactly what such a theorem, combined with the
normalization of $\kapr$, would supply.
\end{remark}

\begin{remark}[Main conjectures at additive primes $2$ and $3$]\label{rmk:closure:t-bad-gap}
For $E/\Q$ with additive reduction at $p\in\{2,3\}$, we know of no
Iwasawa main conjecture giving an equality of characteristic ideals with
explicit local correction factors.  This is to be distinguished from the
signed theories at primes of good supersingular reduction and from other
local constructions.  The hypothesis $\GammaPD$ excludes the prime $2$.  The
hypothesis $\GammaSP$ requires separately supplied signed applicability
and comparison data, and no such comparison at an additive prime $2$ is
established here.
\end{remark}

\begin{remark}[Where the local residuals are carried]\label{rmk:closure:mode-b-residuals}
In the bookkeeping of the Six Birds framework, the difficulties of the
previous two remarks and of the signed local comparison appear as three
residuals: a rank-$r$ leading-term residual, a signed corestriction
residual, and a local descent residual.  They are not additional results.
They are carried, respectively, by the hypotheses $\GammaGZ$, $\GammaSP$
and $\GammaPD$, and naming them does not reduce what has to be supplied.
\end{remark}

\section{The conditional theorem}
\label{sec:landing}

We can now state the main result.  \Cref{fig:landing-chain-diagram}
shows its logical structure: which hypotheses enter the proof of the
scalar identity, and which are carried along.

\begin{figure}[tbp]
\centering
\begin{tikzpicture}[
  font=\small,
  input/.style={draw, rounded corners=2pt, align=center, minimum width=4.7cm,
    minimum height=1.35cm, fill=black!3},
  box/.style={draw, rounded corners=2pt, align=center, minimum width=3.0cm,
    minimum height=1.0cm, fill=black!3},
  result/.style={draw, rounded corners=2pt, align=center, minimum width=4.0cm,
    minimum height=1.35cm, fill=black!8},
  arrow/.style={-{Latex[length=2mm]}, thick}
]
\node[input] (inputs) at (-4.9,0.25) {shell-closure record\\
  three recognition records\\headline imports};
\node[box] (master) at (0,0.25) {closure master\\theorem applies};
\node[result] (landing) at (4.9,0.25) {\(\rhoE=0\)\\
  {\scriptsize readout equivalence}\\Strong-BSD identity};

\draw[arrow] (inputs) -- (master);
\draw[arrow] (master) -- (landing);
\node[align=center, font=\footnotesize] at (0,-1.15)
  {conditional on the supplied records, the shell residual vanishes and the scalar identity follows};
\end{tikzpicture}
\caption{The conditional landing chain from the shell-closure record,
recognition records, and headline imports to the Strong-BSD scalar identity.}
\label{fig:landing-chain-diagram}
\end{figure}


\begin{theorem}[Conditional Strong BSD]\label{thm:closure:strong-bsd-conditional}
Let $\SelBSD$ be a shell with curve $E$ and analytic rank $r$, and let
$C$ be its closure operator.  Assume:
\begin{enumerate}
\item the recognition predicate $\PiBSD(E)$
  (\Cref{def:closure:pi-bsd});
\item the Cassels--Tate input $\chiCTp$ at $E$ and some prime $p$, and the
  Selmer-complex pairing input $\AofE$ at $E$
  (\Cref{def:closure:chi-ct-p-import,def:closure:a-e-import});
\item if $r\ge2$, the archimedean comparison input at $E$, of rank $r$
  (\Cref{def:closure:beilinson-import}).
\end{enumerate}
Then
\[
  \frac{L^{(r)}(E,1)}{r!}
  =
  \frac{\#\Sha(E/\Q)\,\Reg_{\mathrm{NT}}(E)\,\Omega_E\,\Tam(E)}
       {\#E(\Q)_{\tors}^{\,2}},
\]
equivalently $\rhoE=0$, and the asserted components of the inputs in
(2) and (3) hold.%
\footnote{The formalization records this theorem as a conditional
schema, because it returns the components of the inputs in (2) and (3)
by projection.  Its scalar conclusion, however, is derived: Lean proves
it from items (3) and (4) of $\PiBSD(E)$ by the algebra in the proof
below.  A compiled dependency check confirms that the two declarations
isolating this algebra depend on no other declaration of the project,
and that neither the derivation from $\PiBSD(E)$ nor the full theorem
depends on \Cref{thm:closure:chi-ct-p-comparison},
\Cref{thm:closure:oc-eta-formula}, \Cref{thm:closure:composite-signature},
the statements of \Cref{sec:aor-instance}, or the supplied applicability,
computability, falsifiability and ablation fields of the shell.}
\end{theorem}

\begin{proof}
If $r\in\{0,1\}$, item (4) of $\PiBSD(E)$ is the identity $L^*=R_E$.
If $r\ge2$, item (3) supplies a $\GammaGZ$ record whose factors are those
of the shell and whose constant is $\kapr=\#\Sha/d$, with $d$ the squared
torsion order.  Its hypothesis $\psi_-(E,r)=0$ says
$L^*-\kapr\Reg\,\Omega\,\Tam=0$, hence $L^*=\kapr\Reg\,\Omega\,\Tam$ by the
first law of the shell.  Substituting $\kapr=\#\Sha/d$ and applying the
second law $(a/d)\cdot b=(a\cdot b)/d$ three times gives
\[
  L^*
  =\Bigl(\frac{\#\Sha}{d}\Bigr)\Reg\,\Omega\,\Tam
  =\frac{\#\Sha\,\Reg\,\Omega\,\Tam}{d}
  =R_E .
\]
In both cases $\rhoE=0$ by \Cref{thm:closure:pi-bsd-iff-strong-bsd}.  The
remaining assertions are components of the inputs.
\end{proof}

Three comments make the content of the theorem precise.

\emph{What the proof uses.}  The scalar conclusion is derived from the
rank-split part of $\PiBSD(E)$ alone.  The local hypotheses $\GammaPD$
and $\GammaSP$ and the comparison inputs are hypotheses of the theorem,
because they belong to the arithmetic setting in which the recognition
predicate is meant, but the scalar identity does not depend on them.

\emph{How strong the hypothesis is.}  The rank-split scalar premise
contains the identity in rank $0$ and $1$, and expresses it through
normalized fixity in rank at least $2$: $\GammaGZ$ with the normalization
$\kapr=\#\Sha/\#E(\Q)_{\tors}^{\,2}$ is the identity rewritten.  The
remaining hypotheses supply additional records, which the identity does
not reconstruct.  The theorem is therefore not an independent derivation of
Strong BSD from weaker arithmetic statements, and it is not a statement
about minimal hypotheses.  Its content is the exact translation: the
identity is equivalent to the vanishing of $\rhoE$
(\Cref{thm:closure:pi-bsd-iff-strong-bsd}) and, when the scalars form a
field with $2\neq0$, to the fixedness of $x_E$ under the averaging closure
(\Cref{thm:closure:scalar-closure-translation}),
and it follows from the recognition predicate by explicit algebra, with
every arithmetic input named and scoped.

\emph{That the hypotheses are consistent and not vacuous.}  The theorem
has been instantiated on shells with scalars in $\Q$, a label in place of
a curve, all factors equal to $1$, and rank $0$ or $2$.  With leading
coefficient $1$, all hypotheses hold and the conclusion follows; in rank
$0$ no archimedean input is required, and none can be supplied.  With
leading coefficient $2$ and everything else unchanged, the recognition
predicate fails, since it would force $2=1$.  These are formal checks of
the statement, not examples of elliptic curves.

\begin{nonclaim}\label{nonclaim:closure:landing-conditional}
This is not a proof of the Birch--Swinnerton-Dyer conjecture, in its weak
or strong form, for any curve.  The theorem derives the Strong BSD
identity for $E$ from hypotheses whose scalar part contains it in rank at
most one and expresses it through normalized fixity in rank at least two.
We know of no theorem supplying the leading-term identity in rank at
least two, or the main conjectures at additive primes $2$ and $3$
(\Cref{rmk:closure:t-e12-gap,rmk:closure:t-bad-gap}).
\end{nonclaim}

\section{BSD as a non-eliminative AOR instance}
\label{sec:aor-instance}

The conditional landing of \Cref{sec:landing} gives the scalar
Strong-BSD identity under the standing shell-closure record, the named
recognition sources, and the headline imports.
This final technical section records the corresponding
Audited Operational Realisability reading.  In the sense of
\cite{TsiokosAOR2026}, an AOR instance is a refinement-stable,
audit-closed membership predicate: informally, the framework's
``holds in reality'' assertion for an object, meaning that the
residual defects left by refinement have all been accounted for by
approved records.  The reading here is non-eliminative.  The open
arithmetic content is not erased; it is recorded as recognition-source
content, while the mechanical pieces and imported theorem stacks are
accounted for by construction or by approved external citation.

\begin{definition}[BSD AOR-instance object]\label{def:closure:aor-sel-instance}
The AOR-instance object attached to \(\SelBSD\) is the carrier whose
membership witness is the bundle of closure records for the BSD trace
shell.  The bundle has two kinds of entries.  The mechanical entries
are the construction of \(\SelBSD\), the \(\kapr\) normalization, and
the imported theorem stacks recorded through an approved
external-record citation chain.  The substantive entries are the three
recognition-source records of \Cref{sec:recognition-sources}.  The
object is built using the AOR primitives of \cite{TsiokosAOR2026}:
refinement-stable membership, structural-defect accounting, and
discharge records.
\end{definition}

\begin{theorem}[Mechanical-record discharges]\label{thm:closure:aor-mechanical-records}
The mechanical membership components discharge cleanly.  The
\(\SelBSD\) construction and the \(\kapr\) normalization are
discharged by construction.  The \(\chiCTp\), \(\AofE\), and
Beilinson import structures are discharged by an approved
external-record citation chain: they are cited as imported theorem
stacks, not reproved.

For the normalization record, we use \citet{TsiokosBSDApparatus2026}.
Fix the cascade convention
\[
  c_E := \Tam(E)=\prod_{\ell} c_{\ell}(E).
\]
Then the cascade leading-term form
\[
  \frac{L^{(r)}(E,1)}{r!}
  =
  \kapr(E)\,\Omega_E\,\Reg_r(E)\,c_E
\]
agrees with the standard Strong-BSD leading-coefficient form
\[
  \frac{L^{(r)}(E,1)}{r!}
  =
  \Omega_E\,\Reg_r(E)\,
  \frac{\#\Sha(E/\Q)\,\Tam(E)}{\#E(\Q)_{\tors}^{\,2}}
\]
if and only if
\[
  \kapr(E)=
  \frac{\#\Sha(E/\Q)}{\#E(\Q)_{\tors}^{\,2}}.
\]
This is a restatement in cascade coordinates, not a proof of the
Strong-BSD identity.
\footnote{Verified in Lean as
\texttt{SixBirdsBSD.Closure.AORInstance.aorMechanicalRecords};
see App~\ref{app:formalization}.}
\end{theorem}

\begin{proof}
The shell itself is one of the constructed objects of the closure
architecture, so its mechanical entry is accounted for by construction.
The displayed \(\kapr\) equivalence is likewise a constructed
normalization record: once the cascade convention \(c_E=\Tam(E)\) is
fixed, it identifies the scalar that makes the cascade leading-term
form and the standard Strong-BSD leading-coefficient form the same
formula.  It does not assert that the formula is true independently of
the conditional landing.

It remains only to account for the imported theorem stacks.  The
\(\chiCTp\) Cassels--Tate stack, the \(\AofE\) Selmer-complex pairing
stack, and the Beilinson comparison stack enter as named external
records.  They are therefore discharged by an approved citation chain,
not by a derivation inside the closure.  These three citation records,
together with the constructed shell and the \(\kapr\) normalization,
exhaust the mechanical part of the AOR membership witness.
\end{proof}

\begin{theorem}[Recognition-source discharges]\label{thm:closure:aor-recognition-discharge}
The three recognition sources discharge as bridged records.  Here a
bridged record means a recognition residual that is reclassified as
supplied closure content rather than eliminated by a derivation.  Each
\(\Gamma_{\mathrm{BSD}}^{*}\) is a primary structural defect of
source type, carrying its forced secondary defects of role, target,
transport, and limit, and is discharged by such a bridged record.  In
particular, none of the three recognition sources is eliminated.
\footnote{Verified in Lean as
\texttt{SixBirdsBSD.Closure.AORInstance.aorRecognitionDischarge};
see App~\ref{app:formalization}.}
\end{theorem}

\begin{proof}
Each of the three recognition sources is an open arithmetic obligation
that the closure names and consumes but does not derive.  In the audit
presentation, such an obligation is therefore recorded as a
source-type structural defect, together with the secondary information
that specifies its role in the argument, its target factor in the BSD
residual, its transport through the shell, and the limit of what the
closure claims about it.

The discharge is by bridging.  That is, the record says that the
obligation is supplied as recognition content and is available to the
conditional landing; it does not say that the obligation has been
proved away.  This applies to the \(p\)-adic descent source, the
signed-Selmer persistence source, and the higher-rank fixity source.
Thus all three recognition records are accounted for, and all three
remain non-eliminated.
\end{proof}

\Cref{tab:aor-discharge-register} lists the mechanical, imported, and
recognition-source atoms that enter the local AOR membership witness.
\begin{table}[tbp]
\centering
\small
\caption{The eight AOR discharge atoms for the BSD trace shell.  The
recognition-source rows are bridged, not zeroed.}
\label{tab:aor-discharge-register}
\begin{tabularx}{\textwidth}{@{}>{\raggedright\arraybackslash}p{0.31\textwidth}
>{\raggedright\arraybackslash}p{0.24\textwidth}
>{\raggedright\arraybackslash}X@{}}
\toprule
Residual atom & Primary type & Status \\
\midrule
\(\SelBSD\) construction & mechanical & by construction \\
\(\kapr\) normalization & mechanical & by construction, citing the
apparatus paper \\
\(\chiCTp\) import stack & imported theorem stack & approved external
record \\
\(\AofE\) Selmer-complex substrate & imported theorem stack & approved
external record \\
Beilinson rank-\(\geq2\) comparison & imported theorem stack & approved
external record \\
\(\GammaPD\) & recognition source & bridged, not eliminated \\
\(\GammaSP\) & recognition source & bridged, not eliminated \\
\(\GammaGZ\) & recognition source & bridged, not eliminated \\
\bottomrule
\end{tabularx}
\end{table}


\begin{theorem}[AOR-instance membership]\label{thm:closure:aor-instance}
\(\SelBSD\) is refinement-stable audit-closed, hence an AOR instance,
under the declared presentation: its declared audit records discharge
the structural defects that remain under refinement.  This is the
framework's ``Strong BSD holds in reality'' statement, and it is
non-eliminative.  The recognition records are reclassified as bridged
records, not discharged of their open arithmetic.
\footnote{Verified in Lean as
\texttt{SixBirdsBSD.Closure.AORInstance.aorInstance} over a narrowed
representation -- the local syntactic refinement-stable predicate, not
the full Audited Operational Realisability meta-theory; the narrowing
is recorded in App~\ref{app:formalization}.}
\end{theorem}

\begin{proof}
By \Cref{thm:closure:aor-mechanical-records}, the mechanical
components of the BSD trace shell are accounted for: the shell
construction and the \(\kapr\) normalization are discharged by
construction, while the imported theorem stacks are discharged by
approved external citation.  By
\Cref{thm:closure:aor-recognition-discharge}, the three recognition
sources are accounted for as bridged records.

These two groups exhaust the declared structural defects that persist
under refinement of the BSD trace shell.  The mechanical records are
closed by construction or citation; the recognition records are
carried forward as supplied recognition content.  Therefore every
remaining defect is matched by a declared audit record, so the shell is
refinement-stable audit-closed under the declared presentation.

The last clause is essential.  The membership statement does not
convert the recognition sources into classical proofs of the open
arithmetic obligations.  It records that those obligations are present,
named, and bridged in the audit, which is exactly the non-eliminative
AOR reading.
\end{proof}

\begin{remark}[AOR partial status]\label{rmk:closure:aor-partial-status}
The AOR-instance reading does not close the open obligations.  The
rank-\(\geq 2\) generalized Gross--Zagier identity and the additive
\(p\leq 3\) Iwasawa main conjecture remain open in the literature.
What the AOR reading records is that the BSD carrier is audit-closed
under the declared refinement-stable presentation, with the
recognition records reclassified as bridged content rather than
discarded.  Only the local syntactic refinement-stable predicate is
mechanized here, not the full AOR meta-theory.
\end{remark}

\section{Scope and non-claims}\label{sec:scope}
% statements-of-record rows resolved here: none

The closure developed in this paper is deliberately conditional.  The
point of the construction is to isolate exactly which recognition
sources and imported theorem stacks would close the Strong-BSD scalar
identity inside the Six Birds closure discipline, and to record what
is not being claimed.  The following seven non-claims are part of the
statement of the paper.  They fix the intended reading of the theorem,
the role of the recognition sources, the status of the imports, and
the meaning of the final AOR membership assertion.

\begin{itemize}[leftmargin=*]
\item[\textbf{NC-1.}]
\textbf{Not unconditional BSD.}  This is not an unconditional
classical proof of BSD, scalar or strong.  The result is theorem-grade
within the Six Birds closure discipline, conditional on the standing
shell-closure record, the named recognition sources, and the headline
imported theorem stacks.  Outside the framework it reads as the
conditional theorem
\[
  \begin{aligned}
  &\SelBSD\text{-closure}
    \wedge\ \GammaPD \wedge \GammaSP \wedge \GammaGZ\\
  &\quad{}\wedge\ \chiCTp\text{-imports}
    \wedge\ \AofE\text{-imports}
    \wedge\ \text{Beilinson-imports}\\
  &\Longrightarrow \mathrm{Strong\ BSD}.
  \end{aligned}
\]
This is the headline reading of the paper.  The closure is a precise
conditional implication, not a replacement for the missing
unconditional arithmetic theorems.

\item[\textbf{NC-2.}]
\textbf{Recognition sources not derived.}  No recognition source is
derived from framework primitives.  Each \(\Gamma_{\mathrm{BSD}}^{*}\)
is closure content of the formed BSD layer: a typed carrier whose
mathematical payload is supplied as a hypothesis and then consumed by
the shell.  The formalization records these carriers as explicit
structure data, not as axioms and not as hidden constants.  Thus the
three sources are neither proved by the framework nor smuggled into
the conclusion; they are named, typed, and kept visible throughout the
argument.

\item[\textbf{NC-3.}]
\textbf{Rank-\(\geq 2\) Gross--Zagier remains open.}  We do not prove
the rank-\(\geq 2\) generalized Gross--Zagier identity.  The relevant
missing identity is the higher-rank comparison between the leading
\(L\)-value and a Neron--Tate regulator determinant of algebraic
cycles, with the necessary normalizing factors made explicit.  The
no-go result used in this paper is an audit result: in the audited
theorem-grade literature, the needed number-field identity is not
available in the required form.  This is not a non-existence claim.
It says only that the present literature does not supply the theorem
needed to remove the corresponding recognition source.

\item[\textbf{NC-4.}]
\textbf{Additive \(p\leq 3\) Iwasawa main conjecture remains open.}
We do not prove the additive \(p\leq 3\) Iwasawa main conjecture with
the explicit local correction factors and characteristic-ideal
identity required by the closure.  This is again an audit-result
no-go, not a claim that such a theorem cannot exist.  The paper treats
the missing small-additive-prime input as open arithmetic content and
does not derive it from the closure framework.  In particular, the
odd-additive recognition carriers do not by themselves discharge the
\(p=2\) part of this obstruction; any such discharge would have to be a
separate supplied record or a future imported theorem.

\item[\textbf{NC-5.}]
\textbf{Imports not eliminated.}  The \(\chiCTp\) comparison remains
conditional on its imported Cassels--Tate stack.  Likewise, the
\(\AofE\) Selmer-complex pairing substrate and the Beilinson
rank-\(\geq 2\) comparison enter as imported theorem stacks.  These
structures are sourced assumptions for this paper, not results proved
here.  The closure derives consequences from them; it does not
eliminate them, reconstruct them, or claim independent proofs of the
classical arithmetic geometry they package.

\item[\textbf{NC-6.}]
\textbf{Anchor-level evidence is not family-level.}  The
overconvergent \(\eta\)-formula sub-residual is verified at the
anchors \(49\mathrm{a}1@7\), \(36\mathrm{a}1@3\), and
\(121\mathrm{b}1@11\).  Those checks are anchor-level numerical
evidence.  They are not a uniform family-level statement over all
elliptic curves or all additive odd primes in the relevant regime.
The conditional formula in \Cref{sec:eta-formula} remains conditional
on its published-import stack and its named nonvanishing/primitivity
input.

\item[\textbf{NC-7.}]
\textbf{AOR is non-eliminative.}  The AOR-instance reading does not
close the open obligations.  It reclassifies the closure records: the
recognition discharges are bridged records, meaning the corresponding
open arithmetic content is carried as supplied recognition content
rather than eliminated.  The rank-\(\geq 2\) generalized
Gross--Zagier identity and the additive \(p\leq 3\) Iwasawa main
conjecture therefore remain open after the AOR reading.  Moreover,
only the local syntactic refinement-stable predicate is mechanized in
this paper, not the full Audited Operational Realisability
meta-theory.
\end{itemize}

These non-claims also govern the wording of the paper.  We do not say
that the paper proves BSD, proves the additive-prime Iwasawa main
conjecture, or proves the higher-rank Gross--Zagier identity.  The
paper lands Strong BSD as a conditional closure.  The standing shell-closure
record is assumed; recognition sources are supplied and carried as
hypotheses; headline imported theorem stacks are imported and cited; the
AOR reading records a non-eliminative
audit-closed membership statement under those declared records.
\begin{table}[tbp]
\centering
\footnotesize
\caption{Summary of the non-claims of \Cref{sec:scope}.}
\label{tab:nonclaims}
\begin{tabularx}{\textwidth}{@{}>{\raggedright\arraybackslash}p{0.05\textwidth}
>{\raggedright\arraybackslash}p{0.27\textwidth}
>{\raggedright\arraybackslash}X@{}}
\toprule
 & Boundary & Reading \\
\midrule
1 & Not a proof of BSD & Conditional translation; the scalar hypothesis has
the strength of the conclusion. \\
2 & Hypotheses not derived & Recognition hypotheses and comparison inputs
are assumed. \\
3 & No necessity & Normalized \(\GammaGZ\) alone gives the scalar identity
in rank \(\ge2\). \\
4 & No arithmetic formed layer & Intertwining readout and original
fixedness are assumed. \\
5 & No computability & Named computability fields are supplied, not
proved. \\
6 & Open inputs & Rank-\(\ge2\) leading terms; additive \(p\in\{2,3\}\)
main conjectures. \\
7 & \(\eta\)-formula & Conditional certificate only; no evidence
claimed. \\
8 & Register & Statuses under a declared refinement; no arithmetic
proved. \\
\bottomrule
\end{tabularx}
\end{table}


\section{Discussion}\label{sec:discussion}
% statements-of-record rows resolved here: none

The classical results bearing on the Strong BSD formula each control part
of it.  Gross--Zagier and Kolyvagin~\citep{GrossZagier1986,Kolyvagin1989}
relate derivatives, Heegner points, Selmer groups and regulators in rank
at most one; Iwasawa main conjectures~\citep{SkinnerUrban2014} give
$p$-parts of the formula at many primes; Cassels--Tate pairings and
Selmer complexes~\citep{Flach1990,Nekovar2006} organize the finite
arithmetic.  None of them gives the whole formula in general.  This paper
does not change that.  What it does is make one precise statement about
how the formula fits a closure framework, and record exactly which inputs
such a formulation needs.

The precise statement is the fixed-point formulation of
\Cref{thm:closure:scalar-closure-translation}: the identity holds exactly
when the pair of its two sides is fixed by the averaging closure.  It is
elementary, and it is exact; in particular it separates fixedness of the
original pair from the trivial closedness of its image.
\Cref{thm:closure:formed-closure-bridge} says what a closure on a richer
layer of arithmetic data would have to satisfy to return the identity: a
readout intertwining it with the averaging closure, and fixedness of the
original state.  Neither is constructed here, and idempotence alone does
not supply them.

The main theorem should be read with its hypotheses in view.  Its scalar
conclusion follows from a recognition hypothesis that contains the
identity in rank at most one and expresses it through a normalized
leading-term identity in rank at least two.  The other hypotheses, the
local recognition hypotheses and the comparison inputs, are carried with
their scopes but do not enter the derivation.  The value of the
formulation therefore lies in its accounting: what is assumed, where each
assumption would come from, what its scope is, and what is not known.
Several points of that accounting are not obvious and are worth
recording.  The local congruences modulo $\Zp^\times$ cannot see a
Tamagawa number prime to $p$ (\Cref{sec:recognition-sources}).  The
determinant ideal of a Selmer complex does not determine whether a
prescribed generator has a square root (\Cref{ex:closure:1913b1}).  The
proposed $\eta$-formula involves an eigenvalue of slope $1/2$, not the
critical slope (\Cref{sec:eta-formula}).  And sensitivity of the BSD
product to its factors does not make any hypothesis necessary
(\Cref{rmk:closure:sensitivity-not-necessity}).

The companion paper~\citep{TsiokosBSDApparatus2026} studies the same
factors from the other side: which data each factor of the BSD formula
retains, with finite adequacy tests and examples of information loss.
Its normalization of $\kapr$ is what identifies the rank-$\ge2$
recognition hypothesis with the Strong BSD identity here.

The same pattern, a target statement written as the fixed point of a
closure with the deep inputs named as recognition hypotheses and an
audit register recording their status, is the general method of the Six
Birds framework~\citep{TsiokosFoundationsIII2026,TsiokosAOR2026}.  For BSD
it yields a conditional theorem of exactly the strength of its
hypotheses, together with a precise description of the inputs that an
unconditional proof along these lines would require.

\section{Conclusion}\label{sec:conclusion}
% statements-of-record rows resolved here: none

We placed the two sides of the Strong BSD formula in a shell and showed
that the identity is equivalent to the vanishing of their difference and
to the fixedness of the pair under an explicit averaging closure.  We
proved that the identity follows from a rank-split recognition
hypothesis, which contains the identity in rank at most one and
expresses it through a normalized leading-term identity in rank at least
two; the theorem is therefore a conditional translation rather than an
independent derivation.  We stated each further input, the local
recognition hypotheses and the comparisons from the literature, with its
verified scope, and showed by examples where these inputs carry
information that coarser data do not: the square-root behaviour of a
prescribed Selmer-complex generator, and Tamagawa numbers invisible to
congruences modulo units.  The leading-term identity in rank at least two,
and main conjectures at additive primes $2$ and $3$, are inputs that, as
far as we know, no current theorem supplies.

\appendix
\crefalias{section}{appendix}
\section{Lean formalization}\label{app:formalization}
\subsection*{Trust base and methodology}

The Lean development for this paper is deliberately mathlib-free.  It
does not import an external mathematics library, and therefore no
mathlib axioms occur in the dependency closure of the BSD declarations.
Every declaration entered in the manifest with theorem status was
audited with \texttt{\#print axioms}.  The accepted trust base is
Lean's standard foundational base: \texttt{propext},
\texttt{Quot.sound}, and \texttt{Classical.choice}, with
\texttt{funext} and the \texttt{choice} recursor appearing as standard
consequences in the audit output.  No project-local axiom lies behind
any theorem-marked result.  The source tree also contains no active
\texttt{axiom}, \texttt{opaque}, \texttt{constant}, \texttt{sorry}, or
\texttt{admit} declarations in the BSD namespace.  This point matters
for the interpretation of the recognition sources and imports: they
are represented as typed structure carriers, never as Lean axioms.

The formalization mechanizes the framework-structural closure
architecture, not the imported arithmetic geometry.  There are three
main disclosure modes.  First, typed structure carriers are used for
the imported theorem stacks and the recognition sources.  The
Cassels--Tate, Selmer-complex pairing, Beilinson, overconvergent
\(\eta\), and cascade imports are records whose mathematical content is
carried as hypothesis fields.  The three recognition sources are
likewise records: their arithmetic content is carried explicitly and
then consumed by later statements.  In all of these cases the content
is carried, not derived, and it is not an axiom of the formalization.

Second, several results are projection-packaged conditional schemas.
The \(\chiCTp\) comparison, the overconvergent \(\eta\)-formula, the
rank-\(\leq 1\) cascade theorem, the composite signature, and the
headline conditional landing are checked projections over explicit
proof-carrying records, conditional on the relevant imports and
recognition sources.  The formalization checks that the projections and
implication chains use the supplied records in the stated way; it does
not independently derive the underlying arithmetic theorems.  For this
reason the body footnotes say that these claims are tracked as
conditional schemas, not that the formalization derives the arithmetic
comparison unconditionally.

Third, some statements are substantive finite Lean proofs of the
structural layer.  The readout-level equivalence, the
master-theorem-applicability statement, and the AOR mechanical and
recognition discharge theorems are proofs about the finite record
architecture and its bookkeeping.  The AOR-instance membership is
checked over a narrowed representation: the local syntactic
refinement-stable predicate used by this paper, not the full Audited
Operational Realisability meta-theory.  Thus Lean checks the closure
architecture---the residual readout algebra, the implication chains,
and the discharge bookkeeping---while certifying none of the imported
arithmetic geometry or the recognition-source content itself.

In coverage terms, the paper has 28 statements of record.  Of these,
21 are mechanized: 4 substantive structural theorems, 5
projection-packaged conditional schemas, 8 typed definition carriers
(including the shell and predicate mirrors, the typed import carriers,
and the AOR-instance object), 3 recognition-source carriers, and 1
narrowed AOR-instance membership.  The remaining 7 are not mechanized:
the two audit-result no-gos, two remarks, and three non-claims.  The
tables below give the per-label coverage and the corresponding Lean
declaration map.
\subsection*{Wording-discipline table}

Every mechanized statement in the body carries a footnote whose wording
is canonical, not editorial.  The wording encodes whether the result is
a substantive finite Lean proof, a projection-packaged conditional
schema, a narrowed-representation check, a typed structure carrier, or
a recognition-source carrier.  Projection-packaging and narrowing are
therefore visible at the point of disclosure; unmechanized statements
carry no footnote.

\begin{center}
\begin{tabularx}{\textwidth}{@{}>{\raggedright\arraybackslash}p{0.36\textwidth}X@{}}
\toprule
Statement class & Footnote convention \\
\midrule
Theorem, substantive finite proof &
Verified in Lean as \texttt{[decl]}; see App~\ref{app:formalization}. \\
\addlinespace
Theorem, projection-packaged conditional schema &
Tracked in the formalization harness as a conditional schema
\texttt{[decl]} --- a checked projection over an explicit
proof-carrying record (conditional on the imports and recognition
sources), not an independent Lean derivation. \\
\addlinespace
Theorem, narrowed representation &
Verified in Lean as \texttt{[decl]} over a narrowed representation
(the local syntactic refinement-stable predicate); the narrowing is
recorded below. \\
\addlinespace
Definition, typed structure carrier (import) &
Encoded in Lean as a typed structure carrier \texttt{[decl]} bundling
the imported theorem(s) as hypothesis fields; not a Lean axiom. \\
\addlinespace
Definition, recognition-source carrier &
Supplied as a recognition source, encoded in Lean as a typed structure
carrier \texttt{[decl]} over a narrowed representation; its arithmetic
content is a hypothesis field, not a Lean axiom. \\
\addlinespace
Definition, faithful (typed mirror) &
No footnote; the declaration appears only in the per-label coverage
table below. \\
\addlinespace
No-go / remark / non-claim (not mechanized) &
No footnote and no Lean identifier in the body; recorded as a
paper-level statement only. \\
\bottomrule
\end{tabularx}
\end{center}
\subsection*{Standing narrowings}

\paragraph{Projection-packaged conditional schemas.}
Five results in this paper are mechanized as checked projections over
explicit proof-carrying records rather than as independent Lean
derivations: the \(\chiCTp\) Cassels--Tate comparison, the
overconvergent \(\eta\)-formula, the rank-\(\leq 1\) cascade theorem,
the composite signature, and the headline conditional landing.  This
means that the formalization checks the projection or implication
chain from the supplied records to the stated conclusion.  For example,
in the landing it checks that the recognition sources and headline imports
are threaded through the shell-readout theorems so that the residual
vanishes and the Strong-BSD scalar identity follows.

The corresponding bound is equally important.  These projections do
not independently derive the underlying arithmetic-geometry inputs:
the Cassels--Tate comparison inputs, the overconvergent \(L\)-value
inputs, and the rank-\(\leq 1\) Iwasawa/Heegner inputs remain carried
hypotheses.  This is why the body footnotes for these five results say
that they are tracked as conditional schemas and are not independent
Lean derivations.  In particular, the conditional landing is a theorem
with the recognition sources and headline imports as explicit hypothesis
parameters; it is not an unconditional Lean proof of BSD.

\paragraph{Typed structure carriers.}
The five imported theorem stacks and the three recognition sources are
encoded as typed structures whose arithmetic payload is a hypothesis
field.  Three of the imported stacks, \(\chiCTp\), \(\AofE\), and the
Beilinson comparison, are the headline imports used directly in the landing
theorem; the overconvergent \(\eta\) and \(T_{\mathrm{CASCADE}}\) stacks are
supporting conditional imports recorded by their own sections.  For the
imports, this means that the external theorem stacks are carried as
structured records rather than re-derived.  For the recognition sources, it
means that the named arithmetic content is available as a field of the
carrier and is then consumed by the closure architecture.  In both cases
the content is carried, not derived, and it never appears as a Lean axiom.

The recognition-source carriers also use a narrowed structural
representation.  The carrier records the relevant arithmetic assertion
in typed form---for example the \(p\)-adic descent congruence, the
signed-Selmer persistence comparison, or the higher-rank fixity
readout---without formalizing the full underlying \(p\)-adic, Selmer,
or regulator theory inside Lean.  The dimension and role of each
carrier are faithful to the intended arithmetic content, but the
background arithmetic theory is not redeveloped in the formalization.

Neither of these narrowings enlarges the Lean trust base.  They delimit
what the finite mechanization claims: the imported arithmetic and
recognition content are explicit carried hypotheses, while the checked
Lean layer verifies the structural use of those hypotheses.
\begin{table}[tbp]
\centering
\small
\caption{Standing representation notes for the closure formalization.}
\label{tab:representation-notes}
\begin{tabularx}{\textwidth}{@{}>{\raggedright\arraybackslash}p{0.25\textwidth}
>{\raggedright\arraybackslash}p{0.30\textwidth}
>{\raggedright\arraybackslash}X@{}}
\toprule
Representation issue & Where it appears & Disclosure effect \\
\midrule
Import structures & \(\chiCTp\), \(\AofE\), Beilinson, \(\eta\), and
cascade imports & Arithmetic theorem stacks are hypothesis fields in typed
records, not Lean axioms or internal derivations. \\
\addlinespace
Recognition carriers & \(\GammaPD\), \(\GammaSP\), \(\GammaGZ\) & The
arithmetic payload is supplied closure content and is consumed by the shell
as a hypothesis. \\
\addlinespace
Projection-packaged landing & Supporting conditionals and
\Cref{thm:closure:strong-bsd-conditional} & The formalization checks the
conditional projection/implication chain over supplied records, not an
independent arithmetic proof. \\
\addlinespace
Local \(\RefStable\) predicate & AOR-instance theorem & The Lean theorem
checks the paper-local syntactic predicate, not the full AOR meta-theory. \\
\addlinespace
Cross-paper \(\kapr\) import & \(\Cref{sec:aor-instance}\) & The
normalization is reused from the apparatus paper and cited there, not
re-mechanized in the closure axis. \\
\bottomrule
\end{tabularx}
\end{table}

\subsection*{Non-eliminative AOR instance}

The AOR-instance membership theorem is mechanized over a narrowed
representation.  The development defines, for this paper, a local
syntactic refinement-stable predicate for the BSD trace shell; it does
not mechanize the full Audited Operational Realisability meta-theory of
\cite{TsiokosAOR2026}.  The checked statement is therefore the following
bounded claim: with respect to this local predicate, the declared audit
records of the BSD trace shell account for the structural defects that
persist under refinement.  It is not a claim that Lean has formalized the
general AOR meta-theory, and it is not a claim of membership in that
full meta-theoretic sense.  This is why the body footnote for the
AOR-instance theorem carries the narrowed-representation qualifier
rather than a bare ``verified in Lean'' formulation.

The non-eliminative content is also explicit in the formal record.  The
three recognition sources are discharged as bridged records: each is
recorded in the audit as a primary source-type structural defect, with
its forced secondary role, target, transport, and limit defects, and is
then reclassified as carried recognition content.  None of the three is
marked as eliminated.  By contrast, the mechanical components---the
trace-shell construction, the \(\kapr\)-normalization, and the imported
theorem stacks---are discharged either by construction or by an approved
external-record citation chain.

Thus the AOR-instance membership is non-eliminative.  It records that
the open arithmetic content has been named, carried, and bridged inside
the audit presentation, not that the content has been closed.  In
particular, the rank-\(\geq 2\) generalized Gross--Zagier identity and
the additive \(p \leq 3\) Iwasawa main conjecture remain open after the
AOR reading; the membership records their place in the declared
presentation rather than proving them.  This subsection fixes the exact
scope of the formal AOR result: a local, refinement-stable,
non-eliminative audit-closed membership, so that neither the headline
conditional landing nor the AOR reading is mistaken for an unconditional
or eliminative theorem.
\subsection*{Per-label coverage}

The table lists every closure statement of record in document order,
together with its environment kind, mechanization fidelity, and body
location.  Rows marked faithful, projection, or narrowed have a Lean
declaration, mapped in the next subsection; rows marked not mechanized
are paper-level statements and carry no Lean declaration.

\begin{small}
\begin{longtable}{@{}>{\ttfamily}p{0.45\textwidth} l l l@{}}
\caption{Lean coverage of the labelled statements in the inventory, in document order (30 rows).  Labels omit the \texttt{closure:} infix.}\label{tab:formalization-coverage}\\
\toprule \normalfont Label & Kind & Lean coverage & Section \\ \midrule \endfirsthead
\multicolumn{4}{@{}l}{\footnotesize\itshape continued from previous page}\\ \toprule \normalfont Label & Kind & Lean coverage & Section \\ \midrule \endhead
\midrule \multicolumn{4}{r@{}}{\footnotesize\itshape continued on next page}\\ \endfoot
\bottomrule \endlastfoot
def:chi-ct-p-import & \normalfont definition & \normalfont Lean definition & \normalfont \Cref{sec:imports} \\
def:a-e-import & \normalfont definition & \normalfont Lean definition & \normalfont \Cref{sec:imports} \\
def:beilinson-import & \normalfont definition & \normalfont Lean definition & \normalfont \Cref{sec:imports} \\
nonclaim:imports-not-derived & \normalfont nonclaim & \normalfont not formalized & \normalfont \Cref{sec:imports} \\
def:gamma-padic-descent & \normalfont definition & \normalfont Lean definition & \normalfont \Cref{sec:recognition-sources} \\
def:gamma-sha-persistence & \normalfont definition & \normalfont Lean definition & \normalfont \Cref{sec:recognition-sources} \\
def:gamma-higher-gz-fixity & \normalfont definition & \normalfont Lean definition & \normalfont \Cref{sec:recognition-sources} \\
nonclaim:recognition-sources-not-derived & \normalfont nonclaim & \normalfont not formalized & \normalfont \Cref{sec:recognition-sources} \\
def:sel-bsd-shell & \normalfont definition & \normalfont Lean definition & \normalfont \Cref{sec:shell-readout} \\
def:pi-bsd & \normalfont definition & \normalfont Lean definition & \normalfont \Cref{sec:shell-readout} \\
thm:pi-bsd-iff-strong-bsd & \normalfont theorem & \normalfont Lean proves & \normalfont \Cref{sec:shell-readout} \\
thm:scalar-closure-translation & \normalfont theorem & \normalfont Lean proves & \normalfont \Cref{sec:shell-readout} \\
thm:formed-closure-bridge & \normalfont theorem & \normalfont Lean proves & \normalfont \Cref{sec:shell-readout} \\
thm:composite-signature & \normalfont theorem & \normalfont conditional schema & \normalfont \Cref{sec:shell-readout} \\
rmk:sensitivity-not-necessity & \normalfont remark & \normalfont not formalized & \normalfont \Cref{sec:shell-readout} \\
thm:chi-ct-p-comparison & \normalfont theorem & \normalfont conditional schema & \normalfont \Cref{sec:chi-ct-p} \\
def:eta-published-imports & \normalfont definition & \normalfont Lean definition & \normalfont \Cref{sec:eta-formula} \\
thm:oc-eta-formula & \normalfont theorem & \normalfont conditional schema & \normalfont \Cref{sec:eta-formula} \\
def:t-cascade-import & \normalfont definition & \normalfont Lean definition & \normalfont \Cref{sec:t-cascade} \\
thm:t-cascade-rank-le-1 & \normalfont theorem & \normalfont conditional schema & \normalfont \Cref{sec:t-cascade} \\
rmk:t-e12-gap & \normalfont remark & \normalfont not formalized & \normalfont \Cref{sec:obstructions} \\
rmk:t-bad-gap & \normalfont remark & \normalfont not formalized & \normalfont \Cref{sec:obstructions} \\
rmk:mode-b-residuals & \normalfont remark & \normalfont not formalized & \normalfont \Cref{sec:obstructions} \\
thm:strong-bsd-conditional & \normalfont theorem & \normalfont conditional schema & \normalfont \Cref{sec:landing} \\
nonclaim:landing-conditional & \normalfont nonclaim & \normalfont not formalized & \normalfont \Cref{sec:landing} \\
def:aor-sel-instance & \normalfont definition & \normalfont Lean definition & \normalfont \Cref{sec:aor-instance} \\
thm:aor-mechanical-records & \normalfont theorem & \normalfont model version & \normalfont \Cref{sec:aor-instance} \\
thm:aor-recognition-discharge & \normalfont theorem & \normalfont conditional schema & \normalfont \Cref{sec:aor-instance} \\
thm:aor-instance & \normalfont theorem & \normalfont model version & \normalfont \Cref{sec:aor-instance} \\
rmk:aor-partial-status & \normalfont remark & \normalfont not formalized & \normalfont \Cref{sec:aor-instance} \\
\end{longtable}
\end{small}

\subsection*{Lean declaration map}

The 21 mechanized closure statements each correspond to a Lean
declaration, listed here for citation and audit.  Labels elide the
\texttt{closure:} prefix, and declarations elide the
\texttt{SixBirdsBSD.Closure.} namespace; the seven not-mechanized
statements have no declaration and are omitted from this map, with full
coverage recorded in \Cref{tab:formalization-coverage}.  The
\(\kapr\)-normalization reused in the AOR mechanical-records theorem is
imported from \citet{TsiokosBSDApparatus2026};
its own declaration belongs to that paper rather than to the closure
declaration map.

\begin{small}
\begin{longtable}{@{}>{\ttfamily}p{0.44\textwidth} >{\ttfamily}p{0.50\textwidth}@{}}
\caption{Lean declaration map for the 21 mechanized closure statements. Labels elide the \texttt{closure:} prefix; declarations elide the \texttt{SixBirdsBSD.Closure.} namespace.}\label{tab:appd-declmap-closure}\\
\toprule \normalfont Paper label & \normalfont Lean declaration \\ \midrule \endfirsthead
\multicolumn{2}{@{}l}{\footnotesize\itshape continued from previous page}\\ \toprule \normalfont Paper label & \normalfont Lean declaration \\ \midrule \endhead
\midrule \multicolumn{2}{r@{}}{\footnotesize\itshape continued on next page}\\ \endfoot
\bottomrule \endlastfoot
def:\allowbreak chi-\allowbreak ct-\allowbreak p-\allowbreak import & Imports.\allowbreak chi\allowbreak CTp\allowbreak Import \\
def:\allowbreak a-\allowbreak e-\allowbreak import & Imports.\allowbreak a\allowbreak EImport \\
def:\allowbreak beilinson-\allowbreak import & Imports.\allowbreak beilinson\allowbreak Import \\
def:\allowbreak gamma-\allowbreak padic-\allowbreak descent & Recognition\allowbreak Sources.\allowbreak gamma\allowbreak Padic\allowbreak Descent \\
def:\allowbreak gamma-\allowbreak sha-\allowbreak persistence & Recognition\allowbreak Sources.\allowbreak gamma\allowbreak Sha\allowbreak Persistence \\
def:\allowbreak gamma-\allowbreak higher-\allowbreak gz-\allowbreak fixity & Recognition\allowbreak Sources.\allowbreak gamma\allowbreak Higher\allowbreak GZFixity \\
def:\allowbreak sel-\allowbreak bsd-\allowbreak shell & Sel\allowbreak Shell.\allowbreak sel\allowbreak BSDShell \\
def:\allowbreak pi-\allowbreak bsd & Sel\allowbreak Shell.\allowbreak pi\allowbreak BSD \\
thm:\allowbreak pi-\allowbreak bsd-\allowbreak iff-\allowbreak strong-\allowbreak bsd & Sel\allowbreak Shell.\allowbreak pi\allowbreak BSDIff\allowbreak Strong\allowbreak BSD \\
thm:\allowbreak master-\allowbreak theorem-\allowbreak applicability & Sel\allowbreak Shell.\allowbreak master\allowbreak Theorem\allowbreak Applicability \\
thm:\allowbreak composite-\allowbreak signature & Sel\allowbreak Shell.\allowbreak composite\allowbreak Signature \\
thm:\allowbreak chi-\allowbreak ct-\allowbreak p-\allowbreak comparison & Chi\allowbreak CTp.\allowbreak chi\allowbreak CTp\allowbreak Comparison \\
def:\allowbreak eta-\allowbreak published-\allowbreak imports & Eta\allowbreak Formula.\allowbreak eta\allowbreak Published\allowbreak Imports \\
thm:\allowbreak oc-\allowbreak eta-\allowbreak formula & Eta\allowbreak Formula.\allowbreak oc\allowbreak Eta\allowbreak Formula \\
def:\allowbreak t-\allowbreak cascade-\allowbreak import & TCascade.\allowbreak t\allowbreak Cascade\allowbreak Import \\
thm:\allowbreak t-\allowbreak cascade-\allowbreak rank-\allowbreak le-\allowbreak 1 & TCascade.\allowbreak t\allowbreak Cascade\allowbreak Rank\allowbreak Le\allowbreak One \\
thm:\allowbreak strong-\allowbreak bsd-\allowbreak conditional & Landing.\allowbreak strong\allowbreak BSDConditional \\
def:\allowbreak aor-\allowbreak sel-\allowbreak instance & AORInstance.\allowbreak aor\allowbreak Sel\allowbreak Instance \\
thm:\allowbreak aor-\allowbreak mechanical-\allowbreak records & AORInstance.\allowbreak aor\allowbreak Mechanical\allowbreak Records \\
thm:\allowbreak aor-\allowbreak recognition-\allowbreak discharge & AORInstance.\allowbreak aor\allowbreak Recognition\allowbreak Discharge \\
thm:\allowbreak aor-\allowbreak instance & AORInstance.\allowbreak aor\allowbreak Instance \\
\end{longtable}
\end{small}


% Per-axis bibliography, including the Paper A cross-citation; role-level
% citations are maintained in the prose.
\clearpage
\begingroup
\raggedright
\small
\bibliographystyle{plainnat}
\bibliography{references}
\endgroup
\end{document}
